\documentclass[11pt,a4paper]{article}
\usepackage{pkpaper}
% generated by scripts/measure.py; do not edit
\newcommand{\pkItems}{44}
\newcommand{\pkItemsStd}{16}
\newcommand{\pkItemsAuthor}{28}
\newcommand{\pkItemsRewrite}{38}
\newcommand{\pkItemsConv}{4}
\newcommand{\pkItemsGap}{2}
\newcommand{\pkProbes}{420}
\newcommand{\pkProbesUnique}{408}
\newcommand{\pkRules}{32}
\newcommand{\pkRulesR}{21}
\newcommand{\pkRulesRK}{10}
\newcommand{\pkRulesK}{1}
\newcommand{\pkRulesRegular}{16}
\newcommand{\pkRulesSyntactic}{11}
\newcommand{\pkRulesSemantic}{5}
\newcommand{\pkRulesLinted}{19}
\newcommand{\pkLintChecks}{21}
\newcommand{\pkLintDetected}{26}
\newcommand{\pkLintUndetected}{18}
\newcommand{\pkLintPkLongest}{15}
\newcommand{\pkCatalanNineteen}{1\,767\,263\,190}
\newcommand{\pkCatalanTwentyFour}{1\,289\,904\,147\,324}
\newcommand{\pkAttachTen}{58\,786}
\newcommand{\pkAttachEight}{4862}
\newcommand{\pkBracketEleven}{4066}
\newcommand{\pkBracketElevenAll}{16\,796}
\newcommand{\pkScopeSeven}{1066}
\newcommand{\pkScopeSevenAll}{5040}
\newcommand{\pkScopeFour}{14}
\newcommand{\pkScopeFourAll}{24}
\newcommand{\pkOpenList}{256}
\newcommand{\pkParseExpNat}{2,53}
\newcommand{\pkParseExpPk}{1,77}
\newcommand{\pkParseWordsMax}{83}
\newcommand{\pkParseOpsNatMax}{15\,345}
\newcommand{\pkParseOpsPkMax}{1976}
\newcommand{\pkParseRatioMax}{7,8}
\newcommand{\pkParseOpsNatTwenty}{496}
\newcommand{\pkParseOpsPkTwenty}{178}
\newcommand{\pkParseTreesTwenty}{16\,796}
\newcommand{\pkParseRatioTwentyFive}{3,1}
\newcommand{\pkParseTreesTwentyFive}{208\,012}
\newcommand{\pkItemsPkSentences}{88}
\newcommand{\pkTokItems}{40}
\newcommand{\pkTokMedMin}{36}
\newcommand{\pkTokMedMax}{39}
\newcommand{\pkProbesGrayUnique}{136}
\newcommand{\pkClaudeCalls}{272}
\newcommand{\pkClaudeTokMedian}{687}
\newcommand{\pkPrecN}{60}
\newcommand{\pkPrecValid}{57}
\newcommand{\pkPrec}{95}
\newcommand{\pkPrecLo}{86}
\newcommand{\pkPrecHi}{98}
\newcommand{\pkRedInstances}{400}
\newcommand{\pkRedSat}{398}
\newcommand{\pkRedUnsat}{2}
\newcommand{\pkTokUnit}{o200k}
\newcommand{\pkTokTokenizers}{4}
\newcommand{\pkTokFailed}{0}
\newcommand{\pkTokMedian}{37}
\newcommand{\pkTokQone}{13}
\newcommand{\pkTokQthree}{63}
\newcommand{\pkTokMean}{49}
\newcommand{\pkTokMin}{−4}
\newcommand{\pkTokMax}{286}
\newcommand{\pkTokCiLo}{23}
\newcommand{\pkTokCiHi}{53}
\newcommand{\pkTokShareLonger}{90}
\newcommand{\pkTokBreakEven}{27}
\newcommand{\pkTokBreakEvenMean}{33}
\newcommand{\pkWordMedian}{38}
\newcommand{\pkWordMean}{52}
\newcommand{\pkCharMedian}{39}
\newcommand{\pkTokDeltaMedian}{0,37}
\newcommand{\pkTokDeltaMean}{0,49}
\newcommand{\pkLegendTokMax}{37}
\newcommand{\pkProbeModels}{6}
\newcommand{\pkProbeModelsGraded}{4}
\newcommand{\pkProbeLocal}{4}
\newcommand{\pkProbeClaude}{2}
\newcommand{\pkProbeQualified}{5}
\newcommand{\pkProbeFailed}{1}
\newcommand{\pkAccPkMin}{54}
\newcommand{\pkAccPkMax}{99}
\newcommand{\pkAgreeItems}{44}
\newcommand{\pkNatSplit}{73}
\newcommand{\pkPkAllCorrect}{27}
\newcommand{\pkNatSplitCount}{32}
\newcommand{\pkPkAllCorrectCount}{12}
\newcommand{\pkErrPk}{17}
\newcommand{\pkBreakEvenStar}{40}
\newcommand{\pkAnchorMin}{65}
\newcommand{\pkAnchorMax}{92}
\newcommand{\pkAccPkQMin}{70}
\newcommand{\pkAccPkQMax}{99}
\newcommand{\pkSteerQMin}{0,19}
\newcommand{\pkSteerQMax}{0,97}
\newcommand{\pkUndecidedQMin}{0}
\newcommand{\pkUndecidedQMax}{86}
\newcommand{\pkConvPlainMin}{50}
\newcommand{\pkConvPlainMax}{75}
\newcommand{\pkConvLegendMin}{50}
\newcommand{\pkConvLegendMax}{100}
\newcommand{\pkConvOrderAll}{20}
\newcommand{\pkConvOrderPlain}{14}
\newcommand{\pkConvOrderLegend}{15}
\newcommand{\pkConvOrAll}{20}
\newcommand{\pkConvOrPlain}{10}
\newcommand{\pkConvOrLegend}{15}
\newcommand{\pkProbeQualifiedConv}{5}
\newcommand{\pkGapOneReadingOne}{1}
\newcommand{\pkGapOnePkOk}{2}
\newcommand{\pkGapTwoReadingOne}{5}
\newcommand{\pkGapTwoPkOk}{2}
\newcommand{\pkHNatMin}{0,00}
\newcommand{\pkHNatMax}{0,91}
\newcommand{\pkHPkMin}{0,00}
\newcommand{\pkHPkMax}{0,89}
\newcommand{\pkSignWins}{0}
\newcommand{\pkProbeLocalQualified}{3}
\newcommand{\pkClaudeItems}{44}
\newcommand{\pkClaudeSplitCount}{5}
\newcommand{\pkClaudeAgreeCount}{30}
\newcommand{\pkClaudeUnsureCount}{9}
\newcommand{\pkClaudePkBothCount}{34}
\newcommand{\pkAccPkClaudeMin}{91}
\newcommand{\pkAccPkClaudeMax}{99}
\newcommand{\pkSteerClaudeMin}{0,84}
\newcommand{\pkSteerClaudeMax}{0,97}
\newcommand{\pkLocalParamMin}{1,6}
\newcommand{\pkLocalParamMax}{4,8}
\newcommand{\pkAccPkLocalMin}{54}
\newcommand{\pkAccPkLocalMax}{82}
\newcommand{\pkSteerLocalMin}{0,12}
\newcommand{\pkSteerLocalMax}{0,50}
\newcommand{\pkHNatLocalMin}{0,00}
\newcommand{\pkHNatLocalMax}{0,91}
\newcommand{\pkHPkLocalMin}{0,00}
\newcommand{\pkHPkLocalMax}{0,89}
\newcommand{\pkSignWinsLocal}{0}
\newcommand{\pkUnstableNat}{10}
\newcommand{\pkUnstablePk}{1}
\newcommand{\pkCorpusStdCfive}{2}
\newcommand{\pkCorpusActs}{5}
\newcommand{\pkCorpusPk}{3}
\newcommand{\pkCorpusActWords}{219\,801}
\newcommand{\pkCorpusActMin}{43}
\newcommand{\pkCorpusActMax}{71}
\newcommand{\pkCorpusPkMin}{0,0}
\newcommand{\pkCorpusPkMax}{3,5}
\newcommand{\pkCorpusOverMin}{12}
\newcommand{\pkCorpusOverMax}{45}


\title{\textbf{Polszczyzna kontrolowana jako metoda redukcji\\ złożoności obliczeniowej w AI}\\[4pt]
\large Ujęcie logiczno-mnogościowe oraz refleksja filozoficzna\\
w świetle hipotezy Sapira-Whorfa}
\author{Maciej Szulia\\[2pt]\normalsize\texttt{maciej.szulia@gmail.com}}
\date{listopad 2026}

\begin{document}
\maketitle

\begin{abstract}
\noindent
Praca bada tezę, że Polszczyzna Kontrolowana (PK) zmniejsza złożoność obliczeniową pracy
z tekstem. W części formalnej język jest parą: zbiór zdań i relacja interpretacji.
Reguła PK jest operatorem, który usuwa zdania albo usuwa odczytania; operatory są
przemienne. Liczba odczytań zdania jest iloczynem po miejscach wyboru. Praca dowodzi
twierdzenia granicznego. Spójność zbioru wymagań bez przeczeń jest rozstrzygalna w czasie
liniowym. Reguła C5 dopuszcza 1 przeczenie w zdaniu, a 1 przeczenie wystarcza do
NP-zupełności. Część pomiarowa zawiera 5 pomiarów. Zdanie z 8 frazami przyimkowymi ma
\pkAttachEight{} odczytania, a zdanie PK ma 1 odczytanie. PK wydłuża tekst: mediana narzutu
wynosi \pkTokMedian{} \% tokenów. Sonda zadaje \pkProbeModels{} modelom językowym pytania
o sytuacje, w których odczytania zdania się różnią. Modele toru Claude odczytują zdania PK
zgodnie z kluczem w co najmniej \pkAccPkClaudeMin{} \% przypadków. Modele lokalne, które
mają od \pkLocalParamMin{} do \pkLocalParamMax{} mld parametrów włącznie, osiągają od \pkAccPkLocalMin{} \%
do \pkAccPkLocalMax{} \% włącznie. Konwencja niezgodna z nawykiem językowym działa na
model dopiero z legendą. Akty prawne mają od \pkCorpusActMin{} do \pkCorpusActMax{}
trafień sprawdzarki na 1000 słów włącznie. Część filozoficzna opisuje PK jako zestaw
kategorii obowiązkowych i rozdziela hipotezę Sapira-Whorfa dla modelu językowego na
3 tezy. Wniosek: PK zmniejsza koszt wyboru odczytania i koszt wnioskowania, a zwiększa
koszt przejścia tekstu przez model.

\smallskip
\noindent\textbf{Słowa kluczowe:} język kontrolowany; niejednoznaczność; złożoność
obliczeniowa; klauzule Horna; modele językowe; hipoteza Sapira-Whorfa.
\end{abstract}


\section{Wprowadzenie}

Tekst, według którego ktoś działa, ma mieć 1 znaczenie. Polska gramatyka zostawia
czytelnikowi wybór w stałych miejscach: brak rodzajników, swobodny szyk, dopełniacz,
spójniki, granice przedziałów. Zdanie \pkq{Dzieci do 3 lat wchodzą bezpłatnie.} ma
2 odczytania: jedno obejmuje dziecko, które ma mniej niż 3 lata; drugie obejmuje dziecko,
które ma mniej niż 4 lata.

Polszczyzna Kontrolowana, w skrócie PK, zamyka te miejsca wyboru \cite{pk}. PK w wersji 0.1
to zwykła polszczyzna oraz \pkRules{} numerowanych reguł w 5 grupach. Standard ma także
moduł F dla raportów i listę słów kluczowych wymagań. Wzorami PK są angielski standard
ASD-STE100 \cite{ste}, słowa kluczowe RFC 2119 \cite{rfc2119}, notacja EARS \cite{ears}
i polskie \pkq{Zasady techniki prawodawczej} \cite{ztp}.

Czytelnikiem tekstu bywa model językowy. Model wykonuje instrukcję, streszcza
regulamin i sprawdza wymaganie. Praca bada tezę z tytułu: PK zmniejsza złożoność
obliczeniową pracy z tekstem. Praca rozdziela tę tezę na 4 pytania.
\begin{enumerate}[label=(\roman*),leftmargin=2.4em]
  \item Jak zapisać PK w języku logiki i teorii mnogości?
  \item Które koszty obliczeniowe PK zmniejsza, a które koszty PK zwiększa?
  \item Ile wynosi narzut długości PK i jak modele językowe odczytują zdania PK?
  \item Co PK mówi o hipotezie Sapira-Whorfa dla człowieka i dla modelu językowego?
\end{enumerate}

\paragraph{Odpowiedzi w skrócie.}
Język jest parą: zbiór zdań i relacja interpretacji. Reguła PK jest operatorem, który usuwa
zdania albo usuwa odczytania (sekcja \ref{sec:model}). PK zmniejsza koszt 3 mechanizmami:
zwija zbiór odczytań, ogranicza rozmiar zdania i daje zdaniu lokalność
(sekcja \ref{sec:complexity}). PK ma 1 koszt: tekst rośnie. Mediana narzutu wynosi
\pkTokMedian{} \% tokenów (sekcja \ref{sec:results}). Spójność zbioru wymagań bez przeczeń
jest rozstrzygalna w czasie liniowym. Reguła C5 dopuszcza 1 przeczenie w zdaniu, a 1
przeczenie wystarcza do NP-zupełności (twierdzenie \ref{thm:np}). Wynik pomiaru na
modelach językowych podaje sekcja \ref{sec:probe}. Refleksję filozoficzną zawiera
sekcja \ref{sec:philosophy}.

\paragraph{Status wyniku.}
Twierdzenia o parsowaniu, o kodowaniu i o klauzulach Horna są znane; praca przywołuje
źródła. Wkładem pracy są: model mnogościowy reguł PK, klasyfikacja \pkRules{} reguł,
twierdzenie graniczne o 1 przeczeniu, zbiór \pkItems{} pozycji testowych, sprawdzarka reguł
oraz 5 pomiarów. Pomiary mają ograniczony zasięg; sekcja \ref{sec:threats} wymienia
zagrożenia dla trafności. Praca opisuje PK 0.1 i nie ocenia innych standardów.

\paragraph{Konwencje tekstu.}
Autor napisał tę pracę w PK. Legenda: \pkq{lub} znaczy co najmniej 1 z członów;
\pkq{albo} znaczy dokładnie 1 z członów. Część dziesiętną liczby oddziela przecinek.
Cytat, przykład zdania zwykłego i tytuł stoją w cudzysłowie; taki fragment leży poza PK.
Każda liczba w tekście pochodzi ze skryptu \texttt{scripts/measure.py}. Sprawdzarka
z sekcji \ref{sec:check} skontrolowała tekst pracy; wynik podaje tabela \ref{tab:corpus}.

\section{Stan wiedzy}
\label{sec:related}

\paragraph{Języki kontrolowane.}
Język kontrolowany to podzbiór języka naturalnego z ograniczoną gramatyką i ograniczonym
słownictwem. Kuhn opisał 100 języków kontrolowanych opartych na angielskim i zaproponował
klasyfikację PENS \cite{kuhn2014}. Klasyfikacja PENS ocenia 4 cechy: precyzję,
ekspresywność, naturalność i prostotę. Kuhn stwierdza o językach klasy P$^1$, czyli
o językach naturalnych, że automatyczna interpretacja jest problemem \pkq{AI-zupełnym}.
Precyzja języka decyduje więc o koszcie interpretacji.

Języki kontrolowane mają 2 rodowody. Pierwszy rodowód to czytelność dla człowieka:
\pkq{Basic English} Ogdena \cite{ogden1930} i przewodnik AECMA z roku 1986, z którego
powstał standard ASD-STE100 \cite{ste}. Drugi rodowód to interpretacja przez maszynę:
Attempto Controlled English ma gramatykę formalną i tłumaczenie na logikę pierwszego rzędu
\cite{fuchs2008}. PK 0.1 należy do pierwszego rodowodu. Standard PK 0.1 składa się z reguł
zapisanych słowami i nie zawiera gramatyki formalnej. W klasyfikacji PENS autor umieszcza
PK 0.1 w klasie P$^2$: język ma mniej niejednoznaczności niż język naturalny i nie ma
podstawy formalnej.

\paragraph{Kontekst polski.}
Polska tradycja zna reguły redagowania tekstów prawnych \cite{ztp} i ruch prostego języka;
norma ISO 24495-1 opisuje prosty język \cite{iso24495}. Prosty język służy łatwemu
zrozumieniu, a PK służy jednemu znaczeniu. Ziembiński opisuje amfibolię, czyli wieloznaczność
składniową zdania \cite{ziembinski2002}. Przybyszewski zbadał zdania przeczące z kwantyfikatorem
ogólnym \cite{przybyszewski2020}. Kapłański opisał kontrolowaną angielszczyznę
jako interfejs baz wiedzy \cite{kaplanski2011}. Autor nie znalazł opublikowanego standardu
kontrolowanej polszczyzny ogólnego przeznaczenia. Przegląd objął wyszukiwarki naukowe
w listopadzie 2026 r. i ma ograniczony zasięg.

\paragraph{Złożoność a kształt języka.}
Church i Patil pokazali, że liczba rozbiorów zdania z ciągiem fraz przyimkowych rośnie jak
liczby Catalana \cite{church1982}. Earley podał algorytm parsowania gramatyk bezkontekstowych
w czasie rzędu $n^3$, a dla gramatyk jednoznacznych w czasie rzędu $n^2$ \cite{earley1970}.
Knuth zdefiniował gramatyki LR($k$), które parser czyta w czasie liniowym \cite{knuth1965};
Leo rozszerzył algorytm Earleya na czas liniowy dla każdej gramatyki LR($k$) \cite{leo1991}.
Barton wykazał, że parsowanie gramatyk ID/LP, które opisują swobodny szyk, jest NP-zupełne
\cite{barton1985}. Pratt-Hartmann i Third zbadali fragmenty języka angielskiego: problem
spełnialności przechodzi od czasu wielomianowego do nierozstrzygalności, gdy fragment
dostaje kolejne konstrukcje \cite{pratt2004,pratt2006}. Wyniki te pokazują wspólną zasadę:
zakres dopuszczonych konstrukcji wyznacza klasę złożoności.

\paragraph{Niejednoznaczność a modele językowe.}
Zbiór AmbigQA zawiera pytania z więcej niż 1 odpowiedzią \cite{min2020}. Liu i współautorzy
zbudowali zbiór AmbiEnt i stwierdzili, że rozpoznanie niejednoznaczności jest dla modeli
trudne \cite{liu2023}. Farquhar i współautorzy mierzą niepewność modelu entropią liczoną
na poziomie znaczeń \cite{farquhar2024}. Piantadosi, Tily i Gibson bronią
niejednoznaczności: kontekst rozstrzyga znaczenie, więc język zyskuje krótsze formy
\cite{piantadosi2012}. Wynik Piantadosiego wyznacza cenę PK: tekst bez oparcia w kontekście
jest dłuższy.

\section{Model logiczno-mnogościowy}
\label{sec:model}

\subsection{Język jako para}

\begin{definition}
Niech $W$ oznacza zbiór sytuacji. \emph{Znaczenie} to podzbiór zbioru $W$. Zbiór znaczeń to
$M=\mathcal{P}(W)$. \emph{Język} to para $(L,R)$: $L$ jest zbiorem zdań, a
$R\subseteq L\times M$ jest \emph{relacją interpretacji}.
\end{definition}

\begin{definition}
\emph{Zbiór odczytań} zdania $s$ to $\sem{s}=\{m\in M:(s,m)\in R\}$. \emph{Stopień
niejednoznaczności} zdania $s$ to liczba $a(s)=|\sem{s}|$. Zdanie $s$ jest
\emph{jednoznaczne}, gdy $a(s)=1$.
\end{definition}

Język $(L,R)$ jest jednoznaczny dokładnie wtedy, gdy relacja $R$ jest funkcją z $L$ do $M$.
Zwykła polszczyzna jest relacją, która nie jest funkcją. Rysunek \ref{fig:relation}
pokazuje 2 zdania zwykłe z 4 odczytaniami i 4 zdania PK, z których każde ma 1 odczytanie.

\rysunek{relation}{Teoria mnogości: relacja interpretacji zwykłej polszczyzny i funkcja,
do której dąży PK.}{fig:relation}

\subsection{Reguła jako operator}

Standard PK zawiera reguły 2 rodzajów. Reguła pierwszego rodzaju zabrania konstrukcji:
reguła A3 zabrania zaimka, reguła B5 zabrania imiesłowu. Reguła drugiego rodzaju zostawia
zdanie i ustala odczytanie. Według reguły B2 pierwszy rzeczownik jest podmiotem. Według
reguły C1 słowo \pkq{lub} znaczy co najmniej 1 człon.

\begin{definition}
\emph{Ograniczenie} to operator $\rho_\varphi(L,R)=(L\cap\varphi,\;R\cap(\varphi\times M))$
dla zbioru zdań $\varphi$. \emph{Konwencja} to operator $\kappa_K(L,R)=(L,\;R\cap K)$ dla
relacji $K\subseteq L\times M$.
\end{definition}

\begin{proposition}
\label{prop:kernel}
Każde ograniczenie i każda konwencja zmniejsza język, zachowuje zawieranie i jest
idempotentna. Każde 2 operatory są przemienne. Złożenie reguł $r_1,\dots,r_n$ daje język
\[
L'=L\cap\bigcap_i\varphi_i,\qquad R'=R\cap\bigcap_j K_j\cap(L'\times M).
\]
\end{proposition}

\begin{proof}
Oba operatory są przecięciem z ustalonym zbiorem. Przecięcie zbiorów jest przemienne,
łączne i idempotentne.
\end{proof}

Stwierdzenie \ref{prop:kernel} ma 2 skutki praktyczne. Kolejność stosowania reguł PK nie
zmienia wyniku. Sprawdzenie tekstu rozkłada się na niezależne kontrole: 1 kontrola
dla 1 reguły.

Cel standardu ma w modelu 2 warunki. Warunek spójności: $\sem{s}'\neq\emptyset$ dla
każdego $s\in L'$, czyli konwencje zostawiają zdaniu co najmniej 1 odczytanie. Warunek
jednoznaczności: $|\sem{s}'|\le 1$ dla każdego $s\in L'$. Przy obu warunkach relacja $R'$
jest funkcją.

\begin{remark}[status celu]
\label{rem:gap}
Standard PK 0.1 formułuje cel słowami: czytelnik, który zna konwencję, odczytuje każde
zdanie w 1 sposób. Cel jest hipotezą standardu. Numerowane reguły PK 0.1 celu nie
gwarantują. Zdanie \pkq{Operatorowi NIE WOLNO włączyć pompy P1 i pompy P2.} spełnia każdą
regułę numerowaną. Zdanie ma 2 odczytania: zakaz dotyczy każdej z 2 pomp albo zakaz dotyczy
włączenia 2 pomp naraz. Zdanie \pkq{Każdy operator obsługuje dokładnie 1 pompę.} także
spełnia reguły i ma 2 odczytania zasięgu. Standard zamyka takie przypadki krokiem 7
kontroli, czyli testem 2 znaczeń. Test 2 znaczeń wymaga oceny znaczenia przez człowieka.
Pozycje G1 i G2 zbioru testowego mierzą te 2 luki (sekcja \ref{sec:probe}).
\end{remark}

\subsection{Miejsca wyboru}

\begin{definition}
\emph{Miejsce wyboru} zdania $s$ to konstrukcja, która dopuszcza co najmniej 2 odczytania
przy ustalonych odczytaniach reszty zdania. Miejsce wyboru $i$ ma zbiór wariantów $C_i$.
Miejsca wyboru są \emph{niezależne}, gdy każda kombinacja wariantów daje inne odczytanie.
\end{definition}

\begin{theorem}
\label{thm:product}
Jeżeli zdanie $s$ ma $k$ niezależnych miejsc wyboru, to
$a(s)=\prod_{i=1}^{k}|C_i|$ oraz $\log_2 a(s)=\sum_{i=1}^{k}\log_2|C_i|$.
Zamknięcie miejsca wyboru $i$ zmniejsza $\log_2 a(s)$ o $\log_2|C_i|$ bitów.
\end{theorem}

\begin{proof}
Zbiór odczytań jest w bijekcji z iloczynem kartezjańskim $C_1\times\dots\times C_k$.
\end{proof}

Twierdzenie \ref{thm:product} tłumaczy, dlaczego reguły PK działają osobno. Każda reguła
zamyka 1 typ miejsca wyboru i odejmuje stałą liczbę bitów. Miejsca wyboru bywają zależne;
wtedy liczba odczytań jest mniejsza niż iloczyn. Pomiar P1 liczy odczytania dokładnie dla
4 konstrukcji (tabela \ref{tab:counts}, rysunek \ref{fig:growth}).

Według reguły B8 zdanie opisu ma co najwyżej 25 słów. Samo ograniczenie długości nie usuwa
niejednoznaczności. Liczba nawiasowań binarnych ciągu 25 słów to liczba Catalana
$C_{24}=\text{\pkCatalanTwentyFour}$. Zbiór odczytań jest więc skończony, a górna granica liczby odczytań
ma rząd $10^{12}$. Jednoznaczność dają dopiero pozostałe reguły.

\subsection{Strefa szara}

\begin{definition}
\emph{Strefa szara} zdania $s$ to zbiór
$G(s)=\bigcup\sem{s}\setminus\bigcap\sem{s}$.
\end{definition}

Strefa szara zawiera sytuacje, w których odczytania dają różną wartość logiczną. Poza
strefą szarą odczytania są zgodne. Strefa szara jest pusta dokładnie wtedy, gdy każde
2 odczytania zdania są równe; dla zdania jednoznacznego $G(s)=\emptyset$.

\begin{example}
Zdanie \pkq{Dzieci do 3 lat wchodzą bezpłatnie.} ma odczytania $m_1=\{w:\text{wiek}<3\}$
i $m_2=\{w:\text{wiek}<4\}$. Strefa szara to $\{w:3\le\text{wiek}<4\}$. Dziecko, które ma
3 lata i 5 miesięcy, leży w strefie szarej.
\end{example}

Model strefy szarej odpowiada nadwartościowaniu Fine'a \cite{fine1975}: zdanie jest
nadprawdziwe w $\bigcap\sem{s}$ i nadfałszywe poza $\bigcup\sem{s}$. Rysunek
\ref{fig:gray} pokazuje strefy szare 5 pozycji testowych, a rysunek \ref{fig:truth}
pokazuje ten sam obiekt w tablicach prawdy. Pomiar P4 opiera się na strefie szarej: sonda
pyta model o sytuację, w której odczytania się różnią.

\rysunek{gray_zone}{Semantyka: 2 odczytania jako zbiory na osi liczbowej i strefa szara,
w której odczytania się różnią.}{fig:gray}

\rysunek{truth_tables}{Logika: tablice prawdy 2 odczytań dla reguł C2, C1 i C4.}{fig:truth}

\subsection{Moc zbioru odczytań}

Zbiór odczytań ma w modelu 3 możliwe moce.
\begin{itemize}[leftmargin=1.6em]
  \item \textbf{Moc 1.} Zdanie jest jednoznaczne. Taki stan jest celem PK.
  \item \textbf{Moc skończona większa niż 1.} Źródłem jest składnia albo słownik: zasięg
  określenia, zaimek, spójnik, granica przedziału. Lista zakończona skrótem \pkq{itp.}
  oznacza nadzbiór zbioru wymienionych elementów; zdanie nie wskazuje, który nadzbiór. Dla dziedziny o $u$ elementach i listy
  o $k$ elementach zdanie ma $2^{u-k}$ odczytań; dla $u=10$ i $k=2$ zdanie ma \pkOpenList{} odczytań.
  \item \textbf{Moc continuum.} Słowo nieostre, na przykład \pkq{regularnie}, dopuszcza
  każdy próg z przedziału liczb rzeczywistych. Zbiór dopuszczalnych uściśleń jest
  nieprzeliczalny \cite{fine1975}. Reguła A5 każe wybrać 1 próg i zapisać próg liczbą
  z jednostką.
\end{itemize}

Osobną klasę tworzy deiksa. Znaczenie słowa \pkq{jutro} jest funkcją kontekstu; Kaplan
nazywa taką funkcję charakterem wyrażenia \cite{kaplan1989}. Reguła D3 wymaga daty.
Zdanie z datą ma stały charakter: znaczenie zdania jest takie samo w każdym dniu czytania.

\subsection{Klasyfikacja reguł}
\label{sec:rules}

Autor przypisał każdej z \pkRules{} reguł 3 cechy: rodzaj operatora, klasę sprawdzarki
i mechanizm redukcji kosztu. Klasyfikację zawiera plik \texttt{data/rules.csv};
rysunek \ref{fig:rules} pokazuje wynik. Ograniczeniem jest \pkRulesR{} reguł.
Ograniczenie i konwencję łączy \pkRulesRK{} reguł. Samą konwencją jest \pkRulesK{} reguła:
reguła E6, która każe ujawnić konwencje w legendzie. Klasyfikacja jest oceną autora
i nie ma drugiego oceniającego.

\rysunek{rules}{Klasyfikacja reguł PK 0.1 według rodzaju operatora, klasy sprawdzarki
i mechanizmu redukcji kosztu.}{fig:rules}

\section{Złożoność obliczeniowa}
\label{sec:complexity}

\subsection{Zadania czytelnika i mechanizmy PK}

Czytelnik tekstu wykonuje 4 zadania. Zadanie PARSUJ wyznacza zbiór odczytań zdania.
Zadanie WYBIERZ wskazuje odczytanie, które zamierzył autor. Zadanie SPRAWDŹ rozstrzyga,
czy tekst spełnia reguły. Zadanie WNIOSKUJ rozstrzyga, czy zbiór zdań jest spójny i co ze
zbioru wynika. PK zmienia koszt tych zadań 3 mechanizmami.
\begin{itemize}[leftmargin=1.6em]
  \item \textbf{M1, zwinięcie zbioru odczytań.} Reguła zmniejsza liczbę $a(s)$ do wartości 1. Zadanie
  WYBIERZ znika.
  \item \textbf{M2, ograniczenie rozmiaru.} Reguły B8, B4, C2 i C5 ograniczają liczbę słów,
  liczbę rzeczowników w ciągu, mieszanie spójników i liczbę przeczeń. Koszt zadania
  PARSUJ dla 1 zdania ma stałą górną granicę, więc koszt dla tekstu rośnie liniowo
  z liczbą zdań.
  \item \textbf{M3, lokalność.} Reguły A2, A3, D3 i E4 usuwają zaimek, deiksę i odesłanie
  bez numeru. Czytelnik interpretuje zdanie bez przeszukiwania reszty tekstu. Tekst jest
  koniunkcją zdań ze wspólnym słownikiem nazw.
\end{itemize}
Kosztem PK jest długość: nazwa zastępuje zaimek, liczba z jednostką zastępuje słowo
nieostre, a 2 zdania zastępują imiesłów.

\subsection{Zadanie PARSUJ}
\label{sec:parse}

Dla gramatyki bezkontekstowej algorytm Earleya działa w czasie rzędu $n^3$; dla gramatyki
jednoznacznej w czasie rzędu $n^2$ \cite{earley1970}. Dla gramatyki LR($k$) istnieje parser
liniowy \cite{knuth1965,leo1991}. Swobodny szyk podnosi koszt: problem rozpoznawania dla
gramatyk ID/LP jest NP-zupełny \cite{barton1985}. Reguła B2 ustala szyk: podmiot,
orzeczenie, dopełnienie. Reguła B6 stawia określenie przy słowie określanym.

Pomiar P2 porównuje 2 gramatyki modelowe. Gramatyka pierwsza pozwala przyłączyć frazę
przyimkową do czasownika albo do rzeczownika, który stoi wcześniej. Gramatyka druga,
zgodna z regułą B6, przyłącza frazę do rzeczownika bezpośrednio przed frazą. Obie gramatyki
generują te same ciągi słów. Wynik podaje sekcja \ref{sec:resparse}. Gramatyki są modelami
1 konstrukcji i nie opisują całej polszczyzny.

\subsection{Zadanie WYBIERZ}

\begin{theorem}
\label{thm:questions}
Niech $Z$ oznacza zamierzone odczytanie zdania $s$, czyli zmienną losową o rozkładzie $P$
na zbiorze $\sem{s}$. Każda strategia pytań z odpowiedzią \pkq{Tak} albo \pkq{Nie}, która zawsze
wskazuje $Z$, zadaje średnio co najmniej $H(Z)=-\sum_m P(m)\log_2P(m)$ pytań. Istnieje
strategia, która zadaje średnio mniej niż $H(Z)+1$ pytań.
\end{theorem}

\begin{proof}
Strategia pytań jest drzewem binarnym, czyli kodem prefiksowym dla zbioru $\sem{s}$.
Oba oszacowania wynikają z twierdzenia Shannona o kodowaniu bez szumu \cite{shannon1948}.
\end{proof}

Wielkość $H(Z)$ jest brakiem informacji w zdaniu. Czytelnik pobiera brakujące bity
z kontekstu, z pytania do autora albo z własnego domysłu. Dla zdania PK brak wynosi 0 bitów.
Dla zdania z 2 równie prawdopodobnymi odczytaniami brak wynosi 1 bit. Czytelnik, który
zgaduje, myli się wtedy w 50 \% przypadków przy każdej strategii zgadywania.

\subsection{Zadanie WNIOSKUJ}
\label{sec:horn}

Praca definiuje rdzeń zdaniowy PK dla wymagań. Rdzeń obejmuje wzór \pkq{Jeżeli W, to D.}
z reguły C3, zakaz mieszania spójników z reguły C2 i limit przeczeń z reguły C5.

\begin{definition}
Niech $A$ oznacza zbiór zdań prostych twierdzących, czyli \emph{atomów}. \emph{Literał} to
atom albo atom z przeczeniem. Język \PKz{} zawiera zdania 2 postaci:
\begin{center}
\pkq{Jeżeli $l_1$ i \dots{} i $l_k$, to $a$.}\qquad
\pkq{Jeżeli $l_1$ i \dots{} i $l_k$, to konflikt.}
\end{center}
gdzie $k\ge0$, każde $l_i$ jest literałem, a $a$ jest atomem. Zdanie drugiej postaci
zapisuje 2 wymagania dla 1 działania: \pkq{MUSI} oraz \pkq{NIE WOLNO}. Język \PKzj{j}
zawiera zdania, w których poprzednik ma co najwyżej $j$ przeczeń. Zbiór zdań $\Gamma$ jest
\emph{spójny}, gdy co najmniej 1 wartościowanie atomów spełnia każde zdanie zbioru $\Gamma$;
zdanie drugiej postaci jest spełnione, gdy poprzednik jest fałszywy.
\end{definition}

\begin{theorem}
\label{thm:horn}
Spójność zbioru zdań języka \PKzj{0} jest rozstrzygalna w czasie liniowym względem
łącznej długości zdań.
\end{theorem}

\begin{proof}
Zdanie \pkq{Jeżeli $a_1$ i \dots{} i $a_k$, to $a$.} jest klauzulą
$\neg a_1\lor\dots\lor\neg a_k\lor a$. Zdanie z konfliktem jest klauzulą
$\neg a_1\lor\dots\lor\neg a_k$. Każda klauzula ma co najwyżej 1 literał dodatni, czyli
jest klauzulą Horna. Algorytm Dowlinga i Galliera rozstrzyga spełnialność formuły Horna
w czasie liniowym \cite{dowling1984}.
\end{proof}

\begin{theorem}
\label{thm:np}
Spójność zbioru zdań języka \PKzj{1} jest problemem NP-zupełnym.
\end{theorem}

\begin{proof}
Problem należy do klasy NP: świadkiem jest wartościowanie. NP-trudność wynika z redukcji
problemu 3-SAT \cite{cook1971}. Niech $C$ będzie klauzulą o 3 literałach nad atomami
$x$, $y$, $z$. Redukcja zależy od liczby literałów dodatnich w klauzuli $C$.
\begin{center}
\begin{tabular}{ll}
\toprule
Klauzula $C$ & Zdania języka \PKzj{1}\\
\midrule
$\neg x\lor\neg y\lor\neg z$ & Jeżeli $x$ i $y$ i $z$, to konflikt.\\
$\neg x\lor\neg y\lor z$ & Jeżeli $x$ i $y$, to $z$.\\
$\neg x\lor y\lor z$ & Jeżeli $x$ i nie $y$, to $z$.\\
$x\lor y\lor z$ & Jeżeli nie $x$, to $u$. Jeżeli $u$ i nie $y$, to $z$.\\
\bottomrule
\end{tabular}
\end{center}
W 3 pierwszych wierszach zdanie jest równoważne klauzuli. W wierszu czwartym $u$ jest
nowym atomem, a zdania dają klauzule $x\lor u$ oraz $\neg u\lor y\lor z$. Jeżeli
wartościowanie spełnia $x\lor y\lor z$, to rozszerzenie $u=\neg x$ spełnia obie klauzule.
W drugą stronę rezolucja względem $u$ daje $x\lor y\lor z$. Każde zdanie ma co najwyżej
1 przeczenie, a liczba zdań jest co najwyżej 2 razy większa niż liczba klauzul.
\end{proof}

Skrypt sprawdził redukcję na \pkRedInstances{} losowych formułach, z których każda ma
co najwyżej 6 zmiennych: pełne wyliczenie wartościowań dało zgodną odpowiedź w każdym
przypadku. Twierdzenie \ref{thm:np} zgadza się z dychotomią Schaefera \cite{schaefer1978}:
zbiór klauzul Horna z klauzulą $x\lor u$ leży poza każdą klasą łatwą.

\begin{remark}[znaczenie dla standardu]
Reguła C5 dopuszcza 1 przeczenie w zdaniu, a reguła C3 zaleca wzór \pkq{Jeżeli nie W,
to D2.} dla przypadku przeciwnego. Zbiór wymagań zgodny z PK 0.1 jest więc zbiorem zdań
języka \PKzj{1}. Zbiór bez przeczeń w poprzednikach leży po stronie czasu liniowego.
Zbiór z 1 przeczeniem w poprzedniku leży po stronie NP-zupełności. NP-zupełność opisuje przypadek najgorszy; programy
rozstrzygające spełnialność radzą sobie z praktycznymi zbiorami wymagań. Pełna PK 0.1
nie ogranicza ekspresywności polszczyzny, więc wnioskowanie w pełnej PK jest
nierozstrzygalne tak samo jak w logice pierwszego rzędu. Rysunek \ref{fig:classes}
zestawia wyniki z fragmentami języka angielskiego \cite{pratt2006}.
\end{remark}

\rysunek{classes}{Teoria złożoności: położenie problemów związanych z PK wśród klas
złożoności.}{fig:classes}

\subsection{Zadanie SPRAWDŹ}
\label{sec:check}

Reguły PK różnią się kosztem kontroli. Autor wyróżnił 3 klasy sprawdzarek
(rysunek \ref{fig:rules}). Automat skończony nad ciągiem słów rozstrzyga \pkRulesRegular{}
z \pkRules{} reguł: listę słów zakazanych, długość zdania, liczbę przeczeń, zapis granicy.
Analizy składni wymaga \pkRulesSyntactic{} reguł, na przykład szyk z reguły B2 i zasięg
z reguły B6. Oceny znaczenia wymaga \pkRulesSemantic{} reguł, na przykład reguła A1:
1 termin dla 1 pojęcia.

Plik \texttt{src/pklint.py} zawiera sprawdzarkę z \pkLintChecks{} kontrolami. Każda
kontrola jest wyrażeniem regularnym nad słowami 1 zdania, więc sprawdzarka działa
w czasie liniowym. Sprawdzarka wskazuje kandydatów; ocenę kończy człowiek. Na zbiorze
testowym sprawdzarka wykrywa \pkLintDetected{} z \pkItems{} zdań zwykłych i nie zgłasza
żadnego z \pkItemsPkSentences{} zdań PK. Pozostałe \pkLintUndetected{} zdań zwykłych wymaga analizy składni
albo znaczenia.

\subsection{Koszt w modelu językowym}
\label{sec:cost}

Koszt przejścia tekstu przez transformator zależy od liczby tokenów $n$: warstwa uwagi
wykonuje pracę rzędu $n^2$ \cite{vaswani2017}. Rozstrzygnięcie niejednoznaczności tej pracy
nie zmienia. Merrill i Sabharwal wykazali, że transformator o precyzji logarytmicznej
liczy tylko funkcje z klasy TC$^0$ \cite{merrill2023}. Model przenosi więc trudniejsze
obliczenia do kolejnych generowanych tokenów. Zdanie, które wymaga wyboru odczytania,
dokłada modelowi pracy; zdanie jednoznaczne tej pracy nie dokłada. Praca nie mierzy tego
efektu wewnątrz modelu; pomiar P4 mierzy zachowanie modelu na wyjściu.

\begin{theorem}
\label{thm:breakeven}
Niech tekst zwykły ma $n$ tokenów, a tekst PK ma $(1+\delta)n$ tokenów. Niech czytelnik
odczytuje tekst zwykły błędnie z prawdopodobieństwem $p$, a tekst PK z prawdopodobieństwem
$e$. Niech po błędnym odczytaniu zadanie wraca do autora i czytelnik czyta tekst od nowa,
a próby są niezależne. Oczekiwany koszt odczytania w tokenach wynosi $n/(1-p)$ dla tekstu
zwykłego i $(1+\delta)n/(1-e)$ dla tekstu PK. Tekst PK jest tańszy dokładnie wtedy, gdy
\begin{equation}
\label{eq:breakeven}
p>1-\frac{1-e}{1+\delta}.
\end{equation}
\end{theorem}

\begin{proof}
Liczba prób ma rozkład geometryczny o wartości oczekiwanej $1/(1-p)$ albo $1/(1-e)$.
Nierówność \eqref{eq:breakeven} wynika z porównania 2 kosztów.
\end{proof}

Dla $e=0$ warunek \eqref{eq:breakeven} ma postać $p>\delta/(1+\delta)$. Model pomija koszt
wykrycia błędu i koszt szkody; oba koszty przesuwają próg na korzyść PK. Model pomija
także pracę autora. Autor pisze tekst 1 raz, a czytelnicy czytają tekst $k$ razy, więc
dodatkowa praca autora rozkłada się na $k$ odczytów. Model językowy czyta instrukcję
systemową przy każdym wywołaniu.

\section{Metoda pomiarów}
\label{sec:method}

Praca zawiera 5 pomiarów. Pomiary P1 i P2 są obliczeniami na modelu formalnym. Pomiary P3
i P4 korzystają ze zbioru pozycji testowych. Pomiar P5 korzysta z korpusu tekstów.

\subsection{Zbiór pozycji testowych}
\label{sec:items}

Zbiór ma \pkItems{} pozycji (plik \texttt{data/items.csv}, dodatek \ref{app:items}).
Pozycja zawiera 6 części: zdanie zwykłe, opis 2 odczytań, zdanie PK dla odczytania 1,
zdanie PK dla odczytania 2 oraz 3 sondy. Sonda to sytuacja z pytaniem, na które
odpowiedzią jest \pkq{Tak} albo \pkq{Nie}. Pierwsza sonda leży w strefie szarej: odczytanie
1 daje jedną odpowiedź, a odczytanie 2 daje odpowiedź przeciwną. Sondy druga i trzecia są
kotwicami: oba odczytania dają odpowiedź \pkq{Tak} w kotwicy drugiej i odpowiedź \pkq{Nie}
w kotwicy trzeciej. Kotwice sprawdzają, czy model rozumie zadanie.

Zdanie zwykłe w \pkItemsStd{} pozycjach pochodzi z przykładów standardu PK 0.1. Standard
powstał przed tą pracą, więc autor nie dobierał tych zdań pod wynik pomiaru. Pozostałe
\pkItemsAuthor{} pozycji napisał autor. Zbiór ma 3 rodzaje pozycji.
\begin{itemize}[leftmargin=1.6em]
  \item \textbf{Przeredagowanie} (\pkItemsRewrite{} pozycji): zdanie PK jest innym ciągiem
  słów niż zdanie zwykłe.
  \item \textbf{Konwencja} (\pkItemsConv{} pozycje, reguły B2 i C1): zdanie PK dla odczytania
  1 jest tym samym ciągiem słów co zdanie zwykłe. Odczytanie ustala konwencja, więc
  pozycja ma dodatkowy wariant z legendą.
  \item \textbf{Luka} (\pkItemsGap{} pozycje): zdanie zwykłe spełnia numerowane reguły PK
  i ma 2 odczytania (uwaga \ref{rem:gap}).
\end{itemize}
Każde zdanie PK w zbiorze przechodzi sprawdzarkę z sekcji \ref{sec:check}. Zdanie PK
o największej liczbie słów ma \pkLintPkLongest{} słów.

\subsection{Pomiary P1 i P2}

Pomiar P1 liczy odczytania dokładnie. Funkcja \texttt{count\_parses} liczy drzewa rozbioru
zdania z $k$ frazami przyimkowymi. Funkcja \texttt{bracketings} wylicza nawiasowania ciągu
członów połączonych na przemian słowami \pkq{i} oraz \pkq{lub} i liczy różne funkcje
logiczne. Funkcja \texttt{scope\_orders} liczy różne porządki zasięgu kwantyfikatorów
\pkq{każdy} i \pkq{pewien} ustawionych na przemian. Dla co najwyżej 4 kwantyfikatorów test
sprawdza wynik semantycznie: wylicza każdą relację na dziedzinie o 2 elementach.

Pomiar P2 uruchamia parser Earleya na zdaniach o budowie: rzeczownik, czasownik,
rzeczownik oraz $k$ fraz przyimkowych, dla $k$ od 0 do 40 włącznie. Zdanie ma wtedy
od 3 do \pkParseWordsMax{} słów włącznie. Operacją jest 1 próba wpisania stanu do tablicy
parsera. Parser nie zawiera poprawki Leo \cite{leo1991}. Wykładnik wzrostu pochodzi
z regresji liniowej logarytmu liczby operacji względem logarytmu liczby słów, dla zdań,
które mają co najmniej 21 słów.

\subsection{Pomiar P3: długość}

Pomiar P3 liczy długość zdania zwykłego i 2 zdań PK w 6 miarach: słowa, znaki oraz tokeny
\pkTokTokenizers{} tokenizerów. Tokenizery \texttt{o200k\_base} i \texttt{cl100k\_base}
pochodzą z biblioteki \texttt{tiktoken}; tokenizery modeli \texttt{Llama 3.2} i \texttt{Gemma 3} pochodzą
z plików modeli. Tokenizery modeli Bielik wymagają zgody licencyjnej w serwisie
Hugging Face, więc pomiar pomija tokenizery modeli Bielik. Narzut pozycji to iloraz: średnia długość
2 zdań PK podzielona przez długość zdania zwykłego, minus 1. Pomiar obejmuje \pkTokItems{} pozycji:
przeredagowania i luki. Pozycje typu konwencja mają narzut 0 z definicji; kosztem
konwencji jest legenda. Dla mediany narzutu skrypt liczy przedział ufności 95 \% metodą
bootstrap z 10\,000 losowań.

\subsection{Pomiar P4: sonda strefy szarej}
\label{sec:probemethod}

Model językowy dostaje instrukcję systemową: \pkq{Odpowiadasz na pytania o znaczenie
tekstu. Odpowiedz jednym słowem: Tak albo Nie.} Wiadomość zawiera tekst, sondę i powtórzone
polecenie odpowiedzi 1 słowem. Tekstem jest zdanie zwykłe, zdanie PK dla odczytania 1 albo
zdanie PK dla odczytania 2. Pozycja typu konwencja ma także 2 warianty z legendą. Razem
pomiar ma \pkProbes{} sond, czyli \pkProbesUnique{} różnych zapytań na model.

Pomiar ma 2 wykonane tory.
\begin{itemize}[leftmargin=1.6em]
  \item \textbf{Tor lokalny.} Program \texttt{ollama} uruchamia model na procesorze
  komputera autora, z temperaturą 0. Skrypt odczytuje prawdopodobieństwa 20 najbardziej
  prawdopodobnych tokenów na pierwszej pozycji odpowiedzi. Wielkość $P(\text{Tak})$ to
  masa tokenów słowa \pkq{Tak} podzielona przez sumę mas tokenów słów \pkq{Tak}
  i \pkq{Nie}. Tor lokalny obejmuje każdą z \pkProbesUnique{} sond.
  \item \textbf{Tor Claude.} Polecenie \texttt{claude -p} wywołuje modele Haiku i Sonnet
  bez narzędzi, bez plików reguł autora i bez pamięci sesji; mediana zużycia to
  \pkClaudeTokMedian{} tokenów wejściowych na wywołanie. Interfejs nie podaje prawdopodobieństw tokenów. Skrypt zadaje
  więc każdą sondę 2 razy i liczy udział odpowiedzi \pkq{Tak}. Tor Claude obejmuje tylko
  sondy strefy szarej: \pkProbesGrayUnique{} zapytań i \pkClaudeCalls{} wywołania na model.
\end{itemize}
Plan przewidywał tor trzeci, przez serwis OpenRouter. Tor trzeci nie powstał. Zabezpieczenie
komputera autora zablokowało odczyt klucza dostępu, a autor zostawił blokadę na czas
pomiaru.

Niech $q$ oznacza prawdopodobieństwo odczytania 1 w strefie szarej. Skrypt liczy 5 miar.
\begin{itemize}[leftmargin=1.6em]
  \item \textbf{Niepewność} $H_2(q)=-q\log_2q-(1-q)\log_2(1-q)$, w bitach, dla zdania
  zwykłego i średnio dla 2 zdań PK. Miara wymaga prawdopodobieństw, więc dotyczy toru
  lokalnego.
  \item \textbf{Sterowalność}: różnica $q$ dla zdania PK odczytania 1 i $q$ dla zdania PK
  odczytania 2. Wartość 1 oznacza pełne sterowanie; wartość 0 oznacza brak wpływu.
  \item \textbf{Trafność PK}: udział zdań PK, dla których model wybiera odczytanie
  zamierzone, czyli $q>0{,}5$ dla odczytania 1 i $q<0{,}5$ dla odczytania 2. Poziom
  przypadku wynosi 50 \%.
  \item \textbf{Trafność kotwic}: udział kotwic z odpowiedzią zgodną z kluczem.
  \item \textbf{Rozbieżność czytelników}: udział pozycji, w których modele wybierają różne
  odczytania zdania zwykłego.
\end{itemize}
Dla niepewności skrypt wykonuje dokładny test znaków na pozycjach typu przeredagowanie.
Test nie ma poprawki na liczbę modeli.

\subsection{Pomiar P5: korpus}

Korpus zawiera \pkCorpusActs{} aktów prawnych i 3 teksty pisane w PK. Akty prawne pochodzą
z serwisu \texttt{api.sejm.gov.pl}: Konstytucja RP oraz teksty jednolite Kodeksu cywilnego,
Kodeksu pracy, Prawa o ruchu drogowym i \pkq{Zasad techniki prawodawczej}. Akty prawne
nie są przedmiotem prawa autorskiego; repozytorium zapisuje adresy, sumy SHA-256 i wyniki
zliczeń. Teksty PK to: standard PK 0.1, praca \pkq{Rachunek generatorów} tego samego autora
oraz niniejsza praca.

Skrypt usuwa z aktu nagłówki stron i numery jednostek redakcyjnych. Sprawdzarka liczy
trafienia każdej z \pkLintChecks{} kontroli; miarą jest liczba trafień na 1000 słów.
Trafienie wskazuje konstrukcję objętą regułą PK. Trafienie nie dowodzi, że zdanie ma
2 odczytania w kontekście aktu. Autor ocenił ręcznie losową próbę \pkPrecN{} trafień
z aktów prawnych; próba nie obejmuje kontroli długości zdania. Ocenę wykonał 1 oceniający.

\subsection{Powtarzalność}

Polecenie \texttt{make all} uruchamia testy, pomiary P1, P2, P3 i P5, rysunki i budowę
pliku PDF. Surowe odpowiedzi modeli leżą w plikach \texttt{results/probe\_*.jsonl};
skrypt \texttt{scripts/measure.py} przelicza z tych plików każdą liczbę pomiaru P4. Ciągła
integracja nie wywołuje modeli. Wagi modeli lokalnych identyfikuje nazwa z rejestru
programu \texttt{ollama}; modele Claude nie są dostępne bez konta, więc tor Claude
odtwarza tylko posiadacz konta.

\section{Wyniki}
\label{sec:results}

\subsection{P1: liczba odczytań}

Tabela \ref{tab:counts} i rysunek \ref{fig:growth} podają liczbę odczytań dla 3 konstrukcji.
Zdanie z 8 frazami przyimkowymi ma \pkAttachEight{} odczytania. Ciąg 11 członów połączonych
na przemian słowami \pkq{i} oraz \pkq{lub} ma \pkBracketElevenAll{} nawiasowań, które dają
\pkBracketEleven{} różnych funkcji logicznych. Zdanie z 4 kwantyfikatorami ma
\pkScopeFourAll{} porządki zasięgu, w tym \pkScopeFour{} logicznie różnych. Wzrost jest
wykładniczy w każdej z 3 konstrukcji. Zdanie PK ma 1 odczytanie przy każdej liczbie
miejsc wyboru.

\begin{table}[tp]
\centering
\caption{P1: liczba logicznie różnych odczytań zdania. Kolumna 2: $k$ fraz przyimkowych po
dopełnieniu. Kolumna 3: $k$ spójników \pkq{i} oraz \pkq{lub} na przemian. Kolumna 4:
$k+1$ kwantyfikatorów \pkq{każdy} i \pkq{pewien} na przemian.}
\label{tab:counts}
\begin{tabular}{rrrr}
\toprule
$k$ & frazy przyimkowe & spójniki & kwantyfikatory\\
\midrule
\wiersze{data/table_counts.tex}
\bottomrule
\end{tabular}
\end{table}

\rysunek{growth}{Kombinatoryka: wzrost liczby odczytań z liczbą miejsc wyboru.}{fig:growth}

\subsection{P2: operacje parsera}
\label{sec:resparse}

Rysunek \ref{fig:parser} pokazuje liczbę operacji parsera. Dla zdania o 21 słowach parser
wykonuje \pkParseOpsNatTwenty{} operacji z gramatyką niejednoznaczną i \pkParseOpsPkTwenty{}
operacji z gramatyką PK. To samo zdanie ma \pkParseTreesTwenty{} drzew rozbioru
w gramatyce niejednoznacznej i 1 drzewo w gramatyce PK. Dla zdania
o \pkParseWordsMax{} słowach liczby operacji wynoszą \pkParseOpsNatMax{}
oraz \pkParseOpsPkMax{}; iloraz to \pkParseRatioMax.

Dopasowany wykładnik wzrostu wynosi \pkParseExpNat{} dla gramatyki niejednoznacznej
i \pkParseExpPk{} dla gramatyki PK. Oba wykładniki są mniejsze niż granice teoretyczne
3 i 2, bo zdania pomiaru mają co najwyżej \pkParseWordsMax{} słów. Gramatyka PK jest
jednoznaczna i ma rekursję prawostronną. Parser Earleya bez poprawki Leo przetwarza taką
gramatykę w czasie kwadratowym; z poprawką Leo czas jest liniowy \cite{leo1991}. Pomiar
potwierdza kierunek zmiany i nie potwierdza czasu liniowego.

W zakresie reguły B8 przewaga gramatyki PK w liczbie operacji jest umiarkowana: dla zdania
o 25 słowach iloraz wynosi \pkParseRatioTwentyFive. Główną korzyścią w tym zakresie jest
liczba drzew: 1 drzewo zamiast \pkParseTreesTwentyFive{} drzew.

\rysunek{parser}{Teoria złożoności: zmierzona liczba operacji parsera Earleya dla
gramatyki niejednoznacznej i dla gramatyki PK.}{fig:parser}

\subsection{P3: narzut długości}
\label{sec:restokens}

PK wydłuża tekst. Tabela \ref{tab:tokens} i rysunek \ref{fig:tokens} podają narzut
w 6 miarach. W tokenach \pkTokUnit{} mediana narzutu wynosi \pkTokMedian{} \%, a kwartyle
wynoszą \pkTokQone{} \% i \pkTokQthree{} \%. Przedział ufności 95 \% dla mediany to
od \pkTokCiLo{} \% do \pkTokCiHi{} \% włącznie. Średnia narzutu wynosi \pkTokMean{} \%.
Zdanie PK jest dłuższe niż zdanie zwykłe w \pkTokShareLonger{} \% przypadków. Wynik jest
zgodny w 6 miarach: mediany leżą w przedziale od \pkTokMedMin{} \% do \pkTokMedMax{} \% włącznie.

Najwyższy narzut mają pozycje o terminach liczonych w dniach (reguła D4) oraz pozycja
o mieszaniu spójników (reguła C2). Zdanie PK podaje tam dane, których zdanie zwykłe
nie zawiera: dzień numer 1 albo strukturę listy. Narzut bliski 0 mają pozycje, w których
PK zamienia 1 słowo, na przykład słowo \pkq{wszystkie} z przeczeniem na słowo \pkq{żaden}.

Pozycje typu konwencja nie zmieniają zdania. Kosztem jest legenda, którą autor podaje
1 raz w tekście; legendy zbioru mają co najwyżej \pkLegendTokMax{} tokenów.

\begin{table}[tp]
\centering
\caption{P3: narzut długości PK w procentach długości zdania zwykłego; \pkTokItems{} pozycji.
Kolumna końcowa: udział zdań PK dłuższych niż zdanie zwykłe.}
\label{tab:tokens}
\begin{tabular}{lrrrrrr}
\toprule
Miara & mediana & kwartyl 1 & kwartyl 3 & minimum & maksimum & dłuższe (\%)\\
\midrule
\wiersze{data/table_tokens.tex}
\bottomrule
\end{tabular}
\end{table}

\rysunek{tokens}{Lingwistyka korpusowa: długość zdania zwykłego i zdania PK oraz rozkład
narzutu w 6 miarach.}{fig:tokens}

\subsection{P4: sonda strefy szarej}
\label{sec:probe}

Pomiar objął \pkProbeModels{} modeli: \pkProbeLocal{} modele toru lokalnego
i \pkProbeClaude{} modele toru Claude. Modele lokalne mają od \pkLocalParamMin{} do \pkLocalParamMax{} mld parametrów włącznie;
modele Bielik opisuje raport techniczny \cite{ociepa2025}. Sprawozdanie z pomiarów podaje
listę biegów. Tabela \ref{tab:probe} podaje wynik każdego modelu,
a rysunek \ref{fig:probe} pokazuje odpowiedzi dla każdej pozycji.

\begin{table}[tp]
\centering
\caption{P4: wynik sondy. L: tor lokalny, C: tor Claude. $n$: liczba pozycji typu
przeredagowanie. $H$: mediana niepewności w bitach. Znaki: liczba pozycji z niepewnością
wyższą dla zdania zwykłego i liczba pozycji z niepewnością wyższą dla zdań PK; $p$: test
znaków. Sterowalność: średnia. Kotwice: trafność w procentach; próg kontroli wynosi 80 \%.}
\label{tab:probe}
\small
\setlength{\tabcolsep}{4.5pt}
\begin{tabular}{@{}lrrrrrrrrr@{}}
\toprule
 & & \multicolumn{2}{c}{$H$ (bity)} & \multicolumn{2}{c}{test znaków} & & trafność & \multicolumn{2}{c}{kotwice (\%)}\\
Model & $n$ & zwykłe & PK & znaki & $p$ & sterow. & PK (\%) & zwykłe & PK\\
\midrule
\wiersze{data/table_probe.tex}
\bottomrule
\end{tabular}
\end{table}

\paragraph{Kontrola kotwic.}
Trafność kotwic modeli lokalnych wynosi od \pkAnchorMin{} \% do \pkAnchorMax{} \% włącznie.
Kontroli nie przeszło \pkProbeFailed{} z \pkProbeLocal{} modeli lokalnych. Model, który
myli się w kotwicach, odpowiada błędnie także tam, gdzie zdanie ma 1 odczytanie. Wynik
takiego modelu w strefie szarej nie mówi więc nic o PK. Praca podaje ten wynik w tabeli
i pomija wynik takiego modelu w miarach zgodności.

\paragraph{Teza B: efekt rejestru.}
Zdania PK sterują odczytaniem modeli toru Claude. Trafność PK wynosi w tym torze
od \pkAccPkClaudeMin{} \% do \pkAccPkClaudeMax{} \% włącznie, a sterowalność
od \pkSteerClaudeMin{} do \pkSteerClaudeMax{} włącznie. Modele toru lokalnego mają
trafność PK od \pkAccPkLocalMin{} \% do \pkAccPkLocalMax{} \% włącznie i sterowalność
od \pkSteerLocalMin{} do \pkSteerLocalMax{} włącznie. Teza B zachodzi więc dla modeli
toru Claude. Efekt rejestru jest w modelach lokalnych słabszy niż w modelach toru Claude.

Niepewność $H$ w torze lokalnym ma medianę od \pkHNatLocalMin{} do \pkHNatLocalMax{} bita
włącznie dla zdań zwykłych. Dla zdań PK mediana wynosi od \pkHPkLocalMin{} do \pkHPkLocalMax{} bita włącznie. Test znaków
wykazał niższą niepewność zdań PK w \pkSignWinsLocal{} z \pkProbeLocal{} modeli lokalnych.
Pomiar nie potwierdza więc, że PK zmniejsza niepewność modeli lokalnych. W torze Claude
2 próby dały różne odpowiedzi dla \pkUnstableNat{} \% zdań zwykłych i dla \pkUnstablePk{} \%
zdań PK.

\paragraph{Rozbieżność czytelników.}
Miarę liczą modele, które przeszły kontrolę kotwic, oraz modele toru Claude: razem
\pkProbeQualified{} modeli. Modele te wybierają różne odczytania zdania zwykłego
w \pkNatSplitCount{} z \pkAgreeItems{} pozycji, czyli w \pkNatSplit{} \% pozycji. Zdanie
zwykłe dzieli więc czytelników tak, jak przewiduje model: zbiór odczytań ma więcej niż
1 element i czytelnicy wybierają różne elementy. Oba zdania PK odczytuje zgodnie z kluczem
każdy z tych modeli w \pkPkAllCorrectCount{} z \pkAgreeItems{} pozycji.

Para modeli toru Claude daje obraz ostrzejszy. Modele Haiku i Sonnet wybierają różne
odczytania zdania zwykłego w \pkClaudeSplitCount{} z \pkClaudeItems{} pozycji i to samo
odczytanie w \pkClaudeAgreeCount{} pozycjach. W \pkClaudeUnsureCount{} pozycjach co najmniej
1 model dał 2 różne odpowiedzi w 2 próbach. Oba zdania PK odczytują zgodnie z kluczem oba
modele w \pkClaudePkBothCount{} z \pkClaudeItems{} pozycji.

\paragraph{Teza C: konwencja i legenda.}
Tabela \ref{tab:conv} podaje wynik pozycji typu konwencja. Reguła B2 ustala, że pierwszy
rzeczownik jest podmiotem. Bez legendy modele odczytują zgodnie z kluczem \pkConvOrderPlain{}
z \pkConvOrderAll{} zdań tej reguły, a z legendą \pkConvOrderLegend{} zdań. Reguła C1
rozdziela słowa \pkq{lub} i \pkq{albo}. Bez legendy modele odczytują zgodnie z kluczem
\pkConvOrPlain{} z \pkConvOrAll{} zdań tej reguły, a z legendą \pkConvOrLegend{} zdań.

Wynik odróżnia 2 konwencje. Konwencja B2 potwierdza nawyk: szyk podmiot,
orzeczenie, dopełnienie jest w polszczyźnie szykiem domyślnym. Konwencja C1 zmienia nawyk:
zwykła polszczyzna nie rozdziela słów \pkq{lub} i \pkq{albo} konsekwentnie. Konwencja
zgodna z nawykiem działa bez legendy. Konwencja niezgodna z nawykiem wymaga legendy.
Legendy wymaga także zdanie, w którym wiedza o świecie podpowiada inny podmiot. Model
Sonnet odczytał zdanie \pkq{Napięcie wyłącza zabezpieczenie.} wbrew szykowi, a z legendą
zgodnie z szykiem. Wniosek opiera się na \pkItemsConv{} pozycjach i jest opisem przypadków.

\begin{table}[tp]
\centering
\caption{P4: pozycje typu konwencja (\pkItemsConv{} pozycje, reguły B2 i C1). Trafność
w procentach i średnia sterowalność bez legendy oraz z legendą.}
\label{tab:conv}
\small
\begin{tabular}{lrrrr}
\toprule
Model & trafność & trafność & sterowalność & sterowalność\\
 & bez legendy & z legendą & bez legendy & z legendą\\
\midrule
\wiersze{data/table_conv.tex}
\bottomrule
\end{tabular}
\end{table}

\paragraph{Luki standardu.}
Zdanie G1, \pkq{Operatorowi NIE WOLNO włączyć pompy P1 i pompy P2.}, spełnia numerowane
reguły PK. Odczytanie 1, czyli zakaz dla każdej pompy, wybiera \pkGapOneReadingOne{}
z \pkProbeQualified{} modeli; pozostałe modele wybierają zakaz włączenia 2 pomp naraz.
Zdanie G2, \pkq{Każdy operator obsługuje dokładnie 1 pompę.}, odczytuje jako zdanie
o różnych pompach \pkGapTwoReadingOne{} z \pkProbeQualified{} modeli. Oba zdania PK luki G1
odczytuje zgodnie z kluczem \pkGapOnePkOk{} modeli, a oba zdania PK luki G2
\pkGapTwoPkOk{} modeli. Modele są w obu lukach zgodne między sobą, więc luka nie dzieli tych czytelników. Luka
zostaje groźna dla autora: jeżeli autor zdania G1 zamierzył zakaz dla każdej pompy, to
modele odczytały zdanie wbrew zamiarowi. Wynik potwierdza uwagę \ref{rem:gap}: zgodność
z regułami numerowanymi nie gwarantuje 1 odczytania.

\paragraph{Pozycja ze spornym kluczem.}
Pozycja K2 ma zdanie PK \pkq{Zawór Z1 jest zamknięty [założenie].} i pytanie, czy tekst
stwierdza fakt. Klucz podaje odpowiedź \pkq{Tak}. Oba modele toru Claude odpowiedziały
\pkq{Nie}. Odpowiedź modeli jest uzasadniona: znacznik źródła mówi, że zdanie jest
założeniem. Praca zostawia pozycję w wyniku jako błąd modeli, więc trafność toru Claude
jest zaniżona o 1 pozycję. Obserwacja pokazuje przy okazji, że znacznik z modułu F zmienia
odczytanie zdania przez model.

\rysunek{probe}{Teoria informacji: prawdopodobieństwo odczytania 1 w strefie szarej dla
zdania zwykłego i dla 2 zdań PK, osobno dla każdego modelu.}{fig:probe}

\subsection{Próg opłacalności}
\label{sec:breakeven}

Rysunek \ref{fig:breakeven} łączy pomiary P3 i P4 wzorem \eqref{eq:breakeven}. Mediana
narzutu wynosi $\delta=\text{\pkTokDeltaMedian}$. Dla czytelnika, który odczytuje każde
zdanie PK poprawnie, próg wynosi \pkTokBreakEven{} \%. Modele po kontroli odczytują zdania
PK błędnie w \pkErrPk{} \% przypadków; próg wynosi wtedy \pkBreakEvenStar{} \%. PK jest
tańsza w modelu kosztu, gdy błąd odczytania zdania zwykłego przekracza próg.

Zdanie ze strefy szarej i z 2 równie prawdopodobnymi odczytaniami daje błąd 50 \%.
Zdanie poza strefą szarą takiego błędu nie daje, a kontekst rozstrzyga część
niejednoznaczności. Błąd odczytania całego tekstu jest więc niższy niż 50 \%. Praca tej
wielkości nie mierzy. Wzór wskazuje, gdzie PK się opłaca: w zdaniach, od których zależy
działanie, i w tekstach, które czytelnik czyta bez kontekstu.

\rysunek{breakeven}{Model kosztu: próg opłacalności PK w funkcji narzutu długości.}{fig:breakeven}

\subsection{P5: miejsca wyboru w aktach prawnych}
\label{sec:rescorpus}

Tabela \ref{tab:corpus} i rysunek \ref{fig:corpus} podają wynik sprawdzarki na korpusie.
Akty prawne mają od \pkCorpusActMin{} do \pkCorpusActMax{} trafień na 1000 słów włącznie.
Teksty pisane w PK mają od \pkCorpusPkMin{} do \pkCorpusPkMax{} trafień na 1000 słów
włącznie. Korpus aktów liczy \pkCorpusActWords{} słów.

W aktach prawnych przeważają 4 kontrole. Kontrola A3 zgłasza zaimki. Kontrola KW zgłasza
słowa \pkq{może} i \pkq{musi} poza zapisem słów kluczowych. Kontrola B8 zgłasza zdania
dłuższe niż 25 słów. Kontrola A5 zgłasza słowa nieostre. Udział zdań dłuższych niż 25 słów wynosi w aktach
od \pkCorpusOverMin{} \% do \pkCorpusOverMax{} \% włącznie. Podział aktu na zdania jest
przybliżony, bo akt składa się z jednostek redakcyjnych, a nie ze zdań.

Ręczna ocena próby potwierdza trafność wskazań: \pkPrecValid{} z \pkPrecN{} trafień
wskazuje konstrukcję objętą regułą. Precyzja wynosi \pkPrec{} \%, a przedział Wilsona 95 \%
to od \pkPrecLo{} \% do \pkPrecHi{} \% włącznie. Czułości sprawdzarki praca nie mierzy.

Wynik wymaga 3 zastrzeżeń. Po pierwsze, trafienie wskazuje miejsce wyboru, a nie błąd
aktu. Język prawny ma własne konwencje i własne procedury wykładni. Po drugie, nieostrość
tekstu prawnego bywa celowa: słowo \pkq{niezwłocznie} zostawia sądowi ocenę przypadku
(sekcja \ref{sec:limits}). Po trzecie, teksty pisane w PK także mają trafienia. Standard
PK 0.1 zawiera \pkCorpusStdCfive{} zdania z 2 przeczeniami, czyli narusza własną regułę C5.
Część trafień w tekstach PK jest fałszywa: sprawdzarka zgłasza przymiotnik w nazwie
\pkq{małe litery} jako słowo nieostre.

\begin{table}[tp]
\centering
\caption{P5: wynik sprawdzarki na korpusie. Mediana: mediana długości zdania w słowach.
Kolumna 4: udział zdań dłuższych niż 25 słów w procentach. Trafienia: trafienia
\pkLintChecks{} kontroli na 1000 słów. Kolumna końcowa: liczba wystąpień słowa \pkq{się}
na 1000 słów; słowo \pkq{się} jest kandydatem do oceny, a nie naruszeniem.}
\label{tab:corpus}
\small
\begin{tabular}{@{}lrrrrr@{}}
\toprule
Tekst & słowa & mediana & ponad 25 (\%) & trafienia & \pkq{się}\\
\midrule
\wiersze{data/table_corpus.tex}
\bottomrule
\end{tabular}
\end{table}

\rysunek{corpus}{Lingwistyka korpusowa: liczba trafień sprawdzarki na 1000 słów w aktach
prawnych i w tekstach pisanych w PK.}{fig:corpus}

\section{Refleksja filozoficzna w świetle hipotezy Sapira-Whorfa}
\label{sec:philosophy}

\subsection{Hipoteza i stan sporu}

Humboldt pisał, że budowa języka wpływa na rozwój umysłowy ludzkości \cite{humboldt1836}.
Sapir twierdził w roku 1929, że człowiek jest zdany na język społeczności, a obraz
świata opiera się na nawykach językowych grupy \cite{sapir1929}. Whorf sformułował w roku
1940 zasadę relatywizmu językowego: te same dane fizyczne prowadzą obserwatorów do tego
samego obrazu świata tylko przy podobnym zapleczu językowym \cite{whorf1940,whorf1956}.
Nazwa \pkq{hipoteza Sapira-Whorfa} pochodzi od następców obu autorów.

Literatura rozróżnia 2 wersje hipotezy. Wersja mocna głosi, że język wyznacza myślenie.
Wersja słaba głosi, że język wpływa na myślenie. Kay i Kempton wykryli wpływ nazw barw na
ocenę podobieństwa barw \cite{kay1984}. Winawer i współautorzy pokazali, że osoby
mówiące po rosyjsku szybciej rozróżniają 2 odcienie niebieskiego, gdy odcienie mają
w języku rosyjskim różne nazwy \cite{winawer2007}. Tom pod redakcją Gumperza i Levinsona
porządkuje spór \cite{gumperz1996}. Wersja mocna nie ma potwierdzenia. Wersja słaba ma
potwierdzenie w wybranych dziedzinach.

\subsection{Tradycja polska}

Polska filozofia ma własny wariant tezy. Ajdukiewicz ogłosił w roku 1934 radykalny
konwencjonalizm: obraz świata zależy od aparatury pojęciowej, czyli od wyboru języka
\cite{ajdukiewicz1934}. Schaff poświęcił hipotezie Sapira-Whorfa monografię
\pkq{Język a poznanie} \cite{schaff1964}. Bartmiński rozwinął pojęcie językowego obrazu
świata \cite{bartminski2006}. Wierzbicka szuka pojęć prostych, które są wspólne
językom \cite{wierzbicka1996}.

Szkoła lwowsko-warszawska wiązała jasność języka z jasnością myślenia. Twardowski
twierdził, że niejasny styl ma źródło w niejasnym myśleniu autora \cite{twardowski1919}.
Tarski wykazał, że ścisła definicja prawdy wymaga języka o określonej strukturze
\cite{tarski1933}. PK korzysta z obu myśli. Reguły PK działają w kierunku przeciwnym do
diagnozy Twardowskiego: wymóg jasnego zapisu wymusza jasną decyzję. Wynik Tarskiego
wyznacza granicę: PK 0.1 nie ma struktury formalnej, a język \PKz{} ma.

\subsection{Ideał jednego znaczenia}

PK należy do tradycji języka idealnego. Leibniz szukał alfabetu myśli, w którym spór
zamienia się w rachunek \cite{leibniz1666}. Frege zbudował pismo pojęciowe i porównał
stosunek tego pisma do języka potocznego ze stosunkiem mikroskopu do oka
\cite{frege1879}. Wittgenstein pisze w \pkq{Traktacie logiczno-filozoficznym}, że język
potoczny przesłania formę myśli (teza 4.002) i że każdą myśl da się wypowiedzieć jasno
(teza 4.116) \cite{wittgenstein1921}. Teza 5.6 wiąże granice języka z granicami świata.
Carnap dodał zasadę tolerancji: zadaniem logika jest dochodzenie do konwencji, a nie
ustanawianie zakazów \cite{carnap1934}. Quine opisał regimentację, czyli parafrazę zdań
potocznych na notację kanoniczną \cite{quine1960}. Montague odrzucił istotną różnicę
teoretyczną między językiem naturalnym a językiem formalnym \cite{montague1970}. Brown opisał
w roku 1960 język logiczny Loglan \cite{brown1960}; Loglan jest nowym językiem, a PK jest
podzbiorem polszczyzny.

PK realizuje ten ideał skromnie. PK nie zastępuje polszczyzny notacją. PK wybiera
z polszczyzny zdania, które czytelnik odczytuje w 1 sposób, i dokłada konwencje. Zasada
tolerancji Carnapa opisuje status PK trafnie: PK jest konwencją, którą autor przyjmuje
dla 1 rodzaju tekstu.

Tradycja ma także krytyków. Późny Wittgenstein twierdzi, że reguła nie wyznacza własnego
zastosowania; zastosowanie ustala praktyka wspólnoty \cite{wittgenstein1953}. Reguła E6
standardu PK mówi to samo językiem inżyniera: konwencja działa tylko dla czytelnika,
który zna konwencję. Grice umieścił wśród maksym sposobu 2 zalecenia: unikaj
wieloznaczności oraz mów zwięźle \cite{grice1975}. PK rozstrzyga konflikt tych maksym na
korzyść jednoznaczności; pomiar P3 podaje cenę. Piantadosi, Tily i Gibson pokazali, że
niejednoznaczność oszczędza wysiłek, gdy kontekst rozstrzyga znaczenie
\cite{piantadosi2012}. PK jest więc uzasadniona tam, gdzie kontekstu brakuje albo gdzie
błąd odczytania kosztuje więcej niż dłuższy tekst.

\subsection{PK jako zestaw kategorii obowiązkowych}

Jakobson, za Boasem, ujął różnicę między językami jednym zdaniem: \pkq{Languages differ
essentially in what they must convey and not in what they may convey}
\cite{jakobson1959}. Polszczyzna każe podać rodzaj rzeczownika i aspekt czasownika.
Polszczyzna pozwala pominąć wykonawcę, ilość, granicę przedziału i źródło wiedzy.

PK zmienia listę kategorii obowiązkowych dla 1 rejestru. Autor tekstu w PK podaje 5 danych.
\begin{itemize}[leftmargin=1.6em,itemsep=0pt]
  \item Wykonawca działania (reguła B1).
  \item Ilość albo nazwa obiektu (reguła A2).
  \item Strona granicy przedziału (reguła D2).
  \item Chwila w kalendarzu (reguła D3).
  \item Źródło wiedzy (moduł F).
\end{itemize}
Moduł F jest odpowiednikiem ewidencjalności. Aikhenvald opisała języki, w których każde zdanie wskazuje źródło
informacji: mówiący widział zdarzenie, słyszał o zdarzeniu albo wnioskuje
\cite{aikhenvald2004}. PK wprowadza tę kategorię do polszczyzny konwencją, a nie gramatyką.

Slobin nazwał wpływ kategorii obowiązkowych \pkq{myśleniem dla mówienia}: mówiący kieruje
uwagę na cechy, które gramatyka każe wyrazić \cite{slobin1996}. Praca stawia analogiczną
hipotezę dla pisania. Autor, który pisze w PK, rozstrzyga fakty, które zwykła polszczyzna
pozwala zostawić bez rozstrzygnięcia. Gdy autor faktu nie zna, standard każe wstawić
znacznik \pkq{[DO USTALENIA]}. Przykład ze standardu pokazuje skalę zjawiska: 3 zdania
instrukcji dają 3 znaczniki braku. PK ujawnia więc niewiedzę autora przed wykonaniem
instrukcji. Hipoteza jest słabą wersją tezy Sapira-Whorfa dla 1 rejestru. Praca tej
hipotezy nie testuje; test wymaga badania z udziałem ludzi.

\subsection{Hipoteza Sapira-Whorfa a model językowy}
\label{sec:swllm}

Model językowy jest funkcją: ciąg tokenów wyznacza rozkład prawdopodobieństwa kolejnego
tokenu. Praca rozdziela pytanie o hipotezę Sapira-Whorfa dla modelu na 3 tezy o różnym
statusie.

\paragraph{Teza A: zależność od tekstu.}
Zadanie dociera do modelu wyłącznie jako tekst. Odpowiedź modelu zależy więc od
sformułowania zadania. Teza A zachodzi z samej budowy modelu i nie wymaga pomiaru.
Dla człowieka mocna wersja hipotezy jest sporna, bo człowiek poznaje świat także poza
językiem. Model językowy takiej drogi nie ma: granice tekstu wejściowego są granicami
zadania. Bender i Koller ostrzegają przed wnioskiem, że system wyuczony na samej formie
języka ma dostęp do znaczenia \cite{bender2020}. Praca stosuje tę ostrożność: słowo
\pkq{odczytanie} oznacza w pomiarze P4 rozkład odpowiedzi modelu, a nie stan umysłu.

\paragraph{Teza B: efekt rejestru.}
Ta sama treść zapisana w PK i w zwykłej polszczyźnie daje różne rozkłady odpowiedzi.
Teza B jest empiryczna; wynik podaje sekcja \ref{sec:probe}.

\paragraph{Teza C: efekt nawyku.}
Model uczy się na tekstach w zwykłej polszczyźnie, więc nie zna konwencji PK. Konwencja
działa na model dopiero wtedy, gdy tekst zawiera legendę. Teza C przenosi regułę E6 na
model językowy; wynik podaje sekcja \ref{sec:probe}.

Literatura dostarcza 2 obserwacji o języku wejścia. Wendler i współautorzy pokazali, że
w modelach rodziny Llama-2 reprezentacje warstw pośrednich leżą bliżej języka angielskiego
niż języka zapytania \cite{wendler2024}. Kim i współautorzy zmierzyli 26 języków w zadaniach
wyszukiwania w tekście o długości 64 tys. i 128 tys. tokenów \cite{kim2025}. Język polski
zajął miejsce 1, a język angielski miejsce 6. Wynik dotyczy 7 zadań tego testu
i nie dowodzi ogólnej przewagi polszczyzny. Oba wyniki mówią o języku naturalnym jako
całości. Niniejsza praca bada węższe pytanie: rejestr wewnątrz 1 języka.

Hipoteza ma więc dla modelu inny kształt niż dla człowieka. U człowieka spór dotyczy
trwałego wpływu języka ojczystego na poznanie. U modelu wpływ tekstu wejściowego jest
doraźny i całkowity, a ciekawe pytanie brzmi: które cechy tekstu zmieniają rozkład
odpowiedzi.

\subsection{Granice}
\label{sec:limits}

\paragraph{Nowomowa.}
Orwell opisał język zaprojektowany po to, żeby zawęzić zakres myśli
\cite{orwell1949}. Model z sekcji \ref{sec:model} pokazuje różnicę między nowomową a PK.
Nowomowa zmniejsza zbiór znaczeń, które użytkownik potrafi wyrazić, czyli obraz relacji
interpretacji w zbiorze $M$. PK zmniejsza zbiory odczytań pojedynczych zdań. Dla każdego
ostrego odczytania zdania zwykłego standard PK wskazuje osobne zdanie PK; tabela
w dodatku \ref{app:items} zawiera \pkItemsPkSentences{} takich zdań. Nowomowa odbiera
rozróżnienia, a PK wymusza rozróżnienia. PK jest także dobrowolna i obejmuje 1 rejestr.

\paragraph{Nieostrość celowa.}
PK usuwa znaczenia nieostre, a nieostrość bywa potrzebna. Ustawodawca pisze
\pkq{niezwłocznie} albo \pkq{należyta staranność}, żeby zostawić sądowi ocenę przypadku.
Pomiar P5 pokazuje, że akty prawne używają takich słów stale. Standard PK zastrzega, że
umowa w PK nadal wymaga oceny prawnika. PK nie nadaje się także do literatury, reklamy
i rozmowy; standard wyłącza te rodzaje tekstu.

\paragraph{Konwencja wymaga wspólnoty.}
Jednoznaczność zdania PK nie jest własnością samego ciągu słów. Zdanie \pkq{Ciśnienie
zmienia przepływ.} jest tym samym ciągiem słów w PK i poza PK. Jednoznaczność powstaje
dopiero dla czytelnika, który zna regułę B2. PK jest więc umową między autorem
a czytelnikiem. Rysunek \ref{fig:timeline} umieszcza PK na tle 3 tradycji, z których
standard korzysta.

\rysunek{timeline}{Historia idei: wybrane prace 3 tradycji, które łączy PK.}{fig:timeline}

\section{Zagrożenia dla trafności}
\label{sec:threats}

\paragraph{Trafność konstrukcji.}
Sonda wymusza odpowiedź \pkq{Tak} albo \pkq{Nie}. Model nie ma opcji \pkq{zdanie jest
niejednoznaczne}. Pytanie sondy wskazuje strefę szarą, więc sonda mierzy wybór odczytania,
a nie wykrycie niejednoznaczności. Zdania pomiaru stoją bez kontekstu. W tekście kontekst
rozstrzyga część niejednoznaczności, więc pomiar zawyża skalę problemu zdań zwykłych.

\paragraph{Trafność wewnętrzna.}
Ten sam autor napisał standard, pozycje testowe, klucz odpowiedzi i sprawdzarkę.
Autor napisał \pkItemsAuthor{} z \pkItems{} pozycji po zaprojektowaniu pomiaru.
Klasyfikację reguł i ocenę precyzji wykonał 1 oceniający. Praca nie ma oceny językoznawcy.
Autor ustalił próg kontroli kotwic, czyli 80 \%, po obejrzeniu wyniku pierwszego modelu.
Pozycja K2 ma sporny klucz (sekcja \ref{sec:probe}).

\paragraph{Trafność zewnętrzna.}
Modele lokalne mają od \pkLocalParamMin{} do \pkLocalParamMax{} mld parametrów włącznie; tor Claude
obejmuje 2 modele.
Wynik nie przenosi się wprost na inne modele. Pozycje opisują instalacje, terminy
i regulaminy; inne dziedziny wymagają osobnego pomiaru. Praca nie zawiera badania
z udziałem ludzi.

\paragraph{Wnioskowanie statystyczne.}
Pomiar P4 ma \pkItemsRewrite{} pozycji typu przeredagowanie, \pkItemsConv{} pozycje typu
konwencja i \pkItemsGap{} pozycje typu luka. Wynik dla konwencji i dla luk jest opisem
przypadków, a nie estymacją. Test znaków nie ma poprawki na liczbę modeli. Tor Claude
ma 2 próby na sondę.

\paragraph{Model formalny.}
Twierdzenie \ref{thm:product} zakłada niezależne miejsca wyboru. Twierdzenie
\ref{thm:questions} wymaga rozkładu zamierzonych odczytań, którego praca nie zna.
Twierdzenia \ref{thm:horn} i \ref{thm:np} dotyczą języka \PKz, a nie pełnej PK 0.1.
Gramatyki pomiaru P2 modelują 1 konstrukcję. Twierdzenie \ref{thm:breakeven} zakłada
niezależne próby; model językowy z temperaturą 0 powtarza ten sam błąd, więc wzór
opisuje zbiór zadań, a nie powtórzenia 1 zadania.

\paragraph{Sprawdzarka.}
Sprawdzarka jest przybliżona. Precyzja pochodzi z próby \pkPrecN{} trafień; czułość jest
nieznana. Sprawdzarka nie obejmuje reguł, które wymagają składni albo znaczenia.

\section{Wnioski}
\label{sec:conclusions}

Praca odpowiada na 4 pytania z wprowadzenia.

\paragraph{Pytanie (i): model.}
Język jest parą: zbiór zdań i relacja interpretacji. Reguła PK jest ograniczeniem albo
konwencją; obie operacje są przecięciem z ustalonym zbiorem, więc reguły są przemienne.
PK dąży do tego, żeby relacja interpretacji była funkcją. Numerowane reguły PK 0.1 tego
celu nie gwarantują: praca wskazuje 2 zdania zgodne z regułami, które mają 2 odczytania.

\paragraph{Pytanie (ii): koszty.}
Teza z tytułu zachodzi w 3 znaczeniach i nie zachodzi w 1 znaczeniu.
\begin{itemize}[leftmargin=1.6em,itemsep=1pt]
  \item PK zmniejsza zbiór odczytań. Zdanie z 8 frazami przyimkowymi ma \pkAttachEight{}
  odczytania; zdanie PK ma 1 odczytanie. Brak informacji w zdaniu PK wynosi 0 bitów.
  \item PK zmniejsza koszt parsowania w modelu gramatyki: zmierzony wykładnik spada
  z \pkParseExpNat{} do \pkParseExpPk.
  \item PK wyznacza granicę kosztu wnioskowania. Zbiór wymagań bez przeczeń ma kontrolę
  spójności w czasie liniowym. Zbiór wymagań z 1 przeczeniem w zdaniu ma kontrolę spójności
  NP-zupełną.
  \item PK nie zmniejsza kosztu przejścia tekstu przez model językowy. Tekst PK jest
  dłuższy: mediana narzutu wynosi \pkTokMedian{} \% tokenów.
\end{itemize}

\paragraph{Pytanie (iii): pomiar.}
Modele toru Claude odczytują zdania PK zgodnie z kluczem w co najmniej
\pkAccPkClaudeMin{} \% przypadków. Modele lokalne osiągają od \pkAccPkLocalMin{} \%
do \pkAccPkLocalMax{} \% włącznie przy poziomie przypadku 50 \%. Zdanie zwykłe dzieli
modele po kontroli w \pkNatSplit{} \% pozycji. Korzyść z PK zależy więc od czytelnika:
PK pomaga modelowi, który czyta polskie zdanie poprawnie, i nie zastępuje tej zdolności.
Próg opłacalności w modelu kosztu wynosi \pkBreakEvenStar{} \% błędnych odczytań zdania
zwykłego.

\paragraph{Pytanie (iv): hipoteza Sapira-Whorfa.}
PK zmienia listę kategorii, które autor wyraża obowiązkowo: wykonawca, ilość, granica,
chwila, źródło wiedzy. Dla człowieka praca stawia hipotezę: pisanie w PK ujawnia autorowi braki
w wiedzy o opisywanej czynności. Dla modelu językowego praca rozdziela hipotezę na 3 tezy.
Zależność odpowiedzi od tekstu zachodzi z budowy modelu. Efekt rejestru zachodzi dla
modeli toru Claude. Efekt nawyku widać w pozycjach typu konwencja: konwencja niezgodna
z nawykiem działa dopiero z legendą.

\paragraph{Zalecenia dla standardu.}
Wyniki dają 4 zalecenia dla kolejnej wersji PK.
\begin{enumerate}[leftmargin=1.8em,itemsep=1pt]
  \item Dodać regułę o zasięgu zakazu nad spójnikiem \pkq{i} (luka G1).
  \item Dodać regułę o zasięgu 2 kwantyfikatorów w 1 zdaniu (luka G2).
  \item Poprawić \pkCorpusStdCfive{} zdania standardu, które mają 2 przeczenia.
  \item Zalecić legendę przy każdym tekście dla modelu językowego, który zawiera słowo
  \pkq{lub} albo słowo \pkq{albo}.
\end{enumerate}

\paragraph{Dalsze prace.}
Praca zostawia 5 zadań.
\begin{itemize}[leftmargin=1.6em,itemsep=1pt]
  \item Gramatyka formalna dla rdzenia PK.
  \item Sprawdzarka z analizą składni.
  \item Pomiar na modelach o większej liczbie parametrów.
  \item Pomiar na tekstach z kontekstem.
  \item Badanie z udziałem ludzi, które sprawdzi hipotezę o ujawnianiu braków wiedzy.
\end{itemize}


\begin{thebibliography}{99}
\small
\bibitem{aikhenvald2004} A. Y. Aikhenvald. \emph{Evidentiality}. Oxford University Press, Oxford 2004.
\bibitem{ajdukiewicz1934} K. Ajdukiewicz. Das Weltbild und die Begriffsapparatur. \emph{Erkenntnis} 4 (1934), 259--287.
\bibitem{ste} ASD. \emph{ASD-STE100 Simplified Technical English}, Issue 9. AeroSpace, Security and Defence Industries Association of Europe, Bruksela 2025. \url{https://www.asd-ste100.org}
\bibitem{bartminski2006} J. Bartmiński. \emph{Językowe podstawy obrazu świata}. Wydawnictwo UMCS, Lublin 2006.
\bibitem{barton1985} G. E. Barton. On the complexity of ID/LP parsing. \emph{Computational Linguistics} 11(4) (1985), 205--218.
\bibitem{bender2020} E. M. Bender, A. Koller. Climbing towards NLU: on meaning, form, and understanding in the age of data. \emph{Proceedings of the 58th Annual Meeting of the ACL} (2020), 5185--5198.
\bibitem{rfc2119} S. Bradner. \emph{Key words for use in RFCs to indicate requirement levels}. RFC 2119, 1997. \url{https://doi.org/10.17487/RFC2119}
\bibitem{brown1960} J. C. Brown. Loglan. \emph{Scientific American} 202(6) (1960), 53--63.
\bibitem{carnap1934} R. Carnap. \emph{Logische Syntax der Sprache}. Springer, Wiedeń 1934.
\bibitem{church1982} K. Church, R. Patil. Coping with syntactic ambiguity or how to put the block in the box on the table. \emph{American Journal of Computational Linguistics} 8(3--4) (1982), 139--149.
\bibitem{cook1971} S. A. Cook. The complexity of theorem-proving procedures. \emph{Proceedings of the 3rd Annual ACM Symposium on Theory of Computing} (1971), 151--158.
\bibitem{dowling1984} W. F. Dowling, J. H. Gallier. Linear-time algorithms for testing the satisfiability of propositional Horn formulae. \emph{Journal of Logic Programming} 1(3) (1984), 267--284.
\bibitem{earley1970} J. Earley. An efficient context-free parsing algorithm. \emph{Communications of the ACM} 13(2) (1970), 94--102.
\bibitem{farquhar2024} S. Farquhar, J. Kossen, L. Kuhn, Y. Gal. Detecting hallucinations in large language models using semantic entropy. \emph{Nature} 630 (2024), 625--630.
\bibitem{fine1975} K. Fine. Vagueness, truth and logic. \emph{Synthese} 30(3--4) (1975), 265--300.
\bibitem{frege1879} G. Frege. \emph{Begriffsschrift, eine der arithmetischen nachgebildete Formelsprache des reinen Denkens}. Louis Nebert, Halle 1879.
\bibitem{fuchs2008} N. E. Fuchs, K. Kaljurand, T. Kuhn. Attempto Controlled English for knowledge representation. \emph{Reasoning Web 2008}, Lecture Notes in Computer Science 5224, Springer (2008), 104--124.
\bibitem{grice1975} H. P. Grice. Logic and conversation. W: P. Cole, J. Morgan (red.), \emph{Syntax and Semantics 3: Speech Acts}. Academic Press, Nowy Jork 1975, 41--58.
\bibitem{gumperz1996} J. J. Gumperz, S. C. Levinson (red.). \emph{Rethinking Linguistic Relativity}. Cambridge University Press, Cambridge 1996.
\bibitem{humboldt1836} W. von Humboldt. \emph{Über die Verschiedenheit des menschlichen Sprachbaues und ihren Einfluss auf die geistige Entwickelung des Menschengeschlechts}. F. Dümmler, Berlin 1836.
\bibitem{iso24495} ISO. \emph{ISO 24495-1:2023 Plain language. Part 1: Governing principles and guidelines}. Genewa 2023.
\bibitem{jakobson1959} R. Jakobson. On linguistic aspects of translation. W: R. A. Brower (red.), \emph{On Translation}. Harvard University Press, Cambridge 1959, 232--239.
\bibitem{kaplan1989} D. Kaplan. Demonstratives. W: J. Almog, J. Perry, H. Wettstein (red.), \emph{Themes from Kaplan}. Oxford University Press, Nowy Jork 1989, 481--563.
\bibitem{kaplanski2011} P. Kapłański. Controlled English interface for knowledge bases. \emph{Studia Informatica} 32(2A) (2011), 485--494.
\bibitem{kay1984} P. Kay, W. Kempton. What is the Sapir-Whorf hypothesis? \emph{American Anthropologist} 86(1) (1984), 65--79.
\bibitem{kim2025} Y. Kim, J. Russell, M. Karpinska, M. Iyyer. \emph{One ruler to measure them all: benchmarking multilingual long-context language models}. arXiv:2503.01996, 2025.
\bibitem{knuth1965} D. E. Knuth. On the translation of languages from left to right. \emph{Information and Control} 8(6) (1965), 607--639.
\bibitem{kuhn2014} T. Kuhn. A survey and classification of controlled natural languages. \emph{Computational Linguistics} 40(1) (2014), 121--170.
\bibitem{leibniz1666} G. W. Leibniz. \emph{Dissertatio de arte combinatoria}. Lipsk 1666.
\bibitem{leo1991} J. M. I. M. Leo. A general context-free parsing algorithm running in linear time on every LR($k$) grammar without using lookahead. \emph{Theoretical Computer Science} 82(1) (1991), 165--176.
\bibitem{liu2023} A. Liu, Z. Wu, J. Michael, A. Suhr, P. West, A. Koller, S. Swayamdipta, N. A. Smith, Y. Choi. We're afraid language models aren't modeling ambiguity. \emph{Proceedings of EMNLP 2023}, 790--807.
\bibitem{ears} A. Mavin, P. Wilkinson, A. Harwood, M. Novak. Easy Approach to Requirements Syntax (EARS). \emph{Proceedings of the 17th IEEE International Requirements Engineering Conference} (2009), 317--322.
\bibitem{merrill2023} W. Merrill, A. Sabharwal. The parallelism tradeoff: limitations of log-precision transformers. \emph{Transactions of the ACL} 11 (2023), 531--545.
\bibitem{min2020} S. Min, J. Michael, H. Hajishirzi, L. Zettlemoyer. AmbigQA: answering ambiguous open-domain questions. \emph{Proceedings of EMNLP 2020}, 5783--5797.
\bibitem{montague1970} R. Montague. English as a formal language. W: B. Visentini i in. (red.), \emph{Linguaggi nella società e nella tecnica}. Edizioni di Comunità, Mediolan 1970, 189--224.
\bibitem{ociepa2025} K. Ociepa, Ł. Flis, R. Kinas, K. Wróbel, A. Gwoździej. \emph{Bielik v3 Small: technical report}. arXiv:2505.02550, 2025.
\bibitem{ogden1930} C. K. Ogden. \emph{Basic English: A General Introduction with Rules and Grammar}. Kegan Paul, Trench, Trubner, Londyn 1930.
\bibitem{orwell1949} G. Orwell. \emph{Nineteen Eighty-Four}. Secker \& Warburg, Londyn 1949. Dodatek: The principles of Newspeak.
\bibitem{piantadosi2012} S. T. Piantadosi, H. Tily, E. Gibson. The communicative function of ambiguity in language. \emph{Cognition} 122(3) (2012), 280--291.
\bibitem{pratt2004} I. Pratt-Hartmann. Fragments of language. \emph{Journal of Logic, Language and Information} 13(2) (2004), 207--223.
\bibitem{pratt2006} I. Pratt-Hartmann, A. Third. More fragments of language. \emph{Notre Dame Journal of Formal Logic} 47(2) (2006), 151--177.
\bibitem{przybyszewski2020} S. Przybyszewski. O użyciu operatorów kwantyfikacji ogólnej w zdaniach przeczących. \emph{Język Polski} 2020, z. 3, 5--19.
\bibitem{quine1960} W. V. O. Quine. \emph{Word and Object}. MIT Press, Cambridge 1960.
\bibitem{sapir1929} E. Sapir. The status of linguistics as a science. \emph{Language} 5(4) (1929), 207--214.
\bibitem{schaefer1978} T. J. Schaefer. The complexity of satisfiability problems. \emph{Proceedings of the 10th Annual ACM Symposium on Theory of Computing} (1978), 216--226.
\bibitem{schaff1964} A. Schaff. \emph{Język a poznanie}. Państwowe Wydawnictwo Naukowe, Warszawa 1964.
\bibitem{shannon1948} C. E. Shannon. A mathematical theory of communication. \emph{Bell System Technical Journal} 27 (1948), 379--423, 623--656.
\bibitem{slobin1996} D. I. Slobin. From \emph{thought and language} to \emph{thinking for speaking}. W: \cite{gumperz1996}, 70--96.
\bibitem{pk} M. Szulia. \emph{Polszczyzna Kontrolowana (PK) 0.1}. 2026, wersja \texttt{15503a4}. \url{https://github.com/maciejszulia/polszczyzna-kontrolowana}
\bibitem{tarski1933} A. Tarski. \emph{Pojęcie prawdy w językach nauk dedukcyjnych}. Prace Towarzystwa Naukowego Warszawskiego, Wydział III, nr 34, Warszawa 1933.
\bibitem{twardowski1919} K. Twardowski. O jasnym i niejasnym stylu filozoficznym. \emph{Ruch Filozoficzny} 5 (1919).
\bibitem{vaswani2017} A. Vaswani, N. Shazeer, N. Parmar, J. Uszkoreit, L. Jones, A. N. Gomez, Ł. Kaiser, I. Polosukhin. Attention is all you need. \emph{Advances in Neural Information Processing Systems} 30 (2017), 5998--6008.
\bibitem{wendler2024} C. Wendler, V. Veselovsky, G. Monea, R. West. Do llamas work in English? On the latent language of multilingual transformers. \emph{Proceedings of the 62nd Annual Meeting of the ACL} (2024), 15366--15394.
\bibitem{whorf1940} B. L. Whorf. Science and linguistics. \emph{Technology Review} 42(6) (1940), 229--231, 247--248.
\bibitem{whorf1956} B. L. Whorf. \emph{Language, Thought, and Reality: Selected Writings}. Red. J. B. Carroll. MIT Press, Cambridge 1956.
\bibitem{wierzbicka1996} A. Wierzbicka. \emph{Semantics: Primes and Universals}. Oxford University Press, Oxford 1996.
\bibitem{winawer2007} J. Winawer, N. Witthoft, M. C. Frank, L. Wu, A. R. Wade, L. Boroditsky. Russian blues reveal effects of language on color discrimination. \emph{Proceedings of the National Academy of Sciences} 104(19) (2007), 7780--7785.
\bibitem{wittgenstein1921} L. Wittgenstein. Logisch-philosophische Abhandlung. \emph{Annalen der Naturphilosophie} 14 (1921). Wydanie polskie: \emph{Tractatus logico-philosophicus}, przeł. B. Wolniewicz. Państwowe Wydawnictwo Naukowe, Warszawa 1970.
\bibitem{wittgenstein1953} L. Wittgenstein. \emph{Philosophical Investigations}. Przeł. G. E. M. Anscombe. Basil Blackwell, Oksford 1953.
\bibitem{ztp} \emph{Zasady techniki prawodawczej}. Tekst jednolity: Dz.U. 2026 poz. 300.
\bibitem{ziembinski2002} Z. Ziembiński. \emph{Logika praktyczna}. Wyd. 25. Wydawnictwo Naukowe PWN, Warszawa 2002.
\end{thebibliography}

\appendix

\section{Legenda konwencji PK użytych w pracy}
\label{app:legend}

\begin{itemize}[leftmargin=1.6em,itemsep=1pt]
  \item \pkq{Lub} znaczy co najmniej 1 z członów. \pkq{Albo} znaczy dokładnie 1 z członów.
  \item \pkq{Co najmniej N} znaczy $x\ge N$. \pkq{Co najwyżej N} znaczy $x\le N$.
  \pkq{Od A do B włącznie} znaczy $A\le x\le B$.
  \item Pierwszy rzeczownik w zdaniu jest podmiotem.
  \item Słowo \pkq{tylko} stoi bezpośrednio przed słowem ograniczanym.
  \item Znacznik w nawiasie kwadratowym, na przykład \pkq{[założenie]}, podaje źródło faktu.
  \item Słowa \pkq{MUSI}, \pkq{NIE WOLNO} i \pkq{MOŻE} pisane wielkimi literami są
  słowami kluczowymi wymagań.
\end{itemize}

\section{Pozycje testowe}
\label{app:items}

Kolumna R podaje rodzaj pozycji: O oznacza przeredagowanie, K oznacza konwencję, L oznacza
lukę. Kolumna Ź podaje źródło zdania: S oznacza standard PK 0.1, A oznacza autora. Zdania
PK i sondy każdej pozycji zawiera plik \texttt{data/items.csv}.

{\footnotesize
\begin{longtable}{@{}llll>{\raggedright\arraybackslash}p{0.38\textwidth}>{\raggedright\arraybackslash}p{0.35\textwidth}@{}}
\toprule
Nr & Reguła & R & Ź & Zdanie zwykłe & Odczytanie 1; odczytanie 2\\
\midrule
\endhead
\wiersze{data/table_items.tex}
\bottomrule
\end{longtable}
}

\section{Powtarzalność}
\label{app:repro}

Kod, dane i surowe odpowiedzi modeli zawiera repozytorium pracy \cite{pk}, gałąź
\texttt{paper-complexity}.
\begin{itemize}[leftmargin=1.6em,itemsep=1pt]
  \item \texttt{make all}: testy, pomiary, rysunki i pliki PDF.
  \item \texttt{python scripts/measure.py}: pomiary P1, P2, P3 i P5 oraz agregaty pomiaru P4.
  \item \texttt{python scripts/probe.py ollama MODEL}: sonda na modelu lokalnym.
  \item \texttt{python scripts/probe.py claude MODEL -{}-subset gray -{}-samples 2}:\\ sonda toru Claude.
  \item \texttt{python -m src.pklint -{}-tex PLIK}: sprawdzarka reguł PK dla pliku LaTeX.
\end{itemize}

\end{document}
